% =====================================================================
%  Rhizome Neural Operator — Technical Report
%  Build:  latexmk            (config in .latexmkrc; PDF lands in pdf/)
% =====================================================================
\documentclass[10pt,a4paper,twocolumn]{article}

% =====================================================================
%  Preamble: packages, layout, environments, macros
% =====================================================================

% ---- Encoding & fonts ------------------------------------------------
\usepackage[T1]{fontenc}
\usepackage[utf8]{inputenc}
% Utopia text + Fourier math (as in the PASA two-column style); falls back to Latin Modern
\IfFileExists{fourier.sty}{\usepackage{fourier}}{\usepackage{lmodern}}
\usepackage{microtype}

% ---- Layout: two-column journal style --------------------------------
\usepackage[a4paper,top=2.0cm,bottom=2.2cm,left=1.7cm,right=1.7cm,
            columnsep=0.6cm,headsep=0.5cm]{geometry}
\setlength{\parindent}{1em}
\setlength{\parskip}{0pt}
\raggedbottom
\usepackage{textcase}
\usepackage[explicit]{titlesec}
\titleformat{\section}{\normalfont\normalsize\bfseries}{\thesection}{0.8em}{\MakeTextUppercase{#1}}
\titleformat{name=\section,numberless}{\normalfont\normalsize\bfseries}{}{0pt}{\MakeTextUppercase{#1}}
\titleformat{\subsection}{\normalfont\normalsize\bfseries}{\thesubsection}{0.8em}{#1}
\titleformat{\subsubsection}{\normalfont\normalsize\itshape}{\thesubsubsection}{0.8em}{#1}
\titlespacing*{\section}{0pt}{2.0ex plus 0.6ex minus 0.3ex}{1.0ex plus 0.2ex}
\titlespacing*{\subsection}{0pt}{1.6ex plus 0.5ex minus 0.3ex}{0.8ex plus 0.2ex}
\titlespacing*{\paragraph}{0pt}{1.0ex plus 0.3ex minus 0.2ex}{0.6em}
\usepackage{fancyhdr}
\pagestyle{fancy}
\fancyhf{}
\fancyhead[L]{\small\itshape Rhizome Neural Operator -- Technical Report}
\fancyhead[R]{\small\thepage}
\renewcommand{\headrulewidth}{0pt}
\fancypagestyle{plain}{\fancyhf{}\fancyhead[R]{\small\thepage}\renewcommand{\headrulewidth}{0pt}}
% float placement tuned for two columns
\renewcommand{\topfraction}{0.9}
\renewcommand{\bottomfraction}{0.6}
\renewcommand{\textfraction}{0.07}
\renewcommand{\floatpagefraction}{0.75}
\renewcommand{\dbltopfraction}{0.9}
\renewcommand{\dblfloatpagefraction}{0.75}
\setcounter{topnumber}{3}
\setcounter{dbltopnumber}{2}
\setcounter{totalnumber}{4}
\setlength{\textfloatsep}{10pt plus 3pt minus 2pt}
\setlength{\dbltextfloatsep}{12pt plus 3pt minus 2pt}
\setlength{\floatsep}{10pt plus 3pt minus 2pt}

% ---- Math ------------------------------------------------------------
\usepackage{amsmath,amssymb,amsthm,mathtools}
\usepackage{bm}
\usepackage{siunitx}

% ---- Figures, tables, algorithms ------------------------------------
\usepackage{graphicx}
\graphicspath{{Figures/}}
\usepackage[font=small,labelfont=bf,labelsep=period,justification=justified,
            singlelinecheck=false,skip=5pt]{caption}
\usepackage{subcaption}
\usepackage{booktabs}
\usepackage{multirow}
\usepackage{float}
\usepackage{placeins}
\usepackage{needspace}
\usepackage{algorithm}
\usepackage{algpseudocode}
\usepackage{tikz}
\usetikzlibrary{arrows.meta,positioning,calc,fit,shapes.geometric}

% ---- Colours, references --------------------------------------------
\usepackage[dvipsnames]{xcolor}
\usepackage[numbers,sort&compress]{natbib}
\usepackage[colorlinks=true,linkcolor=NavyBlue,citecolor=ForestGreen,
            urlcolor=Maroon]{hyperref}
\usepackage[capitalise,noabbrev]{cleveref}

% ---- Theorem-like environments --------------------------------------
\theoremstyle{definition}
\newtheorem{definition}{Definition}[section]
\newtheorem{assumption}[definition]{Assumption}
\newtheorem{example}[definition]{Example}
\theoremstyle{plain}
\newtheorem{proposition}[definition]{Proposition}
\newtheorem{lemma}[definition]{Lemma}
\theoremstyle{remark}
\newtheorem{remark}[definition]{Remark}

% ---- Draft notes -----------------------------------------------------
\newif\ifdraft
\newcommand{\TODO}[1]{\ifdraft{\color{BrickRed}\textbf{[TODO:} #1\textbf{]}}\fi}
\newcommand{\NOTE}[1]{\ifdraft{\color{RoyalBlue}\textit{[Note: #1]}}\fi}

% Figure that falls back to a framed placeholder if the file is missing
% Usage: \figplaceholder[width]{file-in-Figures}{placeholder text}
\newcommand{\figplaceholder}[3][0.8\linewidth]{%
  \IfFileExists{Figures/#2}{\includegraphics[width=#1]{#2}}{%
    \fbox{\parbox[c][0.5\linewidth][c]{\dimexpr#1-2\fboxsep-2\fboxrule}{\centering\color{gray}%
      \textit{Figure placeholder:}\\ \texttt{\detokenize{#2}}\\[0.5em] #3}}}}

% ---- Number sets & generic math -------------------------------------
\newcommand{\R}{\mathbb{R}}
\newcommand{\C}{\mathbb{C}}
\newcommand{\Z}{\mathbb{Z}}
\newcommand{\T}{\mathbb{T}}   
\newcommand{\V}{\mathcal{V}}                  % torus / periodic box
\newcommand{\dd}{\,\mathrm{d}}
\newcommand{\E}{\mathbb{E}}
\DeclarePairedDelimiter{\abs}{\lvert}{\rvert}
\DeclarePairedDelimiter{\norm}{\lVert}{\rVert}
\DeclarePairedDelimiter{\inner}{\langle}{\rangle}
\DeclareMathOperator{\diag}{diag}
\DeclareMathOperator{\softplus}{softplus}
\DeclareMathOperator*{\argmin}{arg\,min}

% ---- Operator learning ----------------------------------------------
\newcommand{\G}{\mathcal{G}}                    % true solution operator
\newcommand{\Gtheta}{\mathcal{G}_\theta}        % learned operator
\newcommand{\Kop}{\mathcal{K}}                  % kernel integral operator
\newcommand{\A}{\mathcal{A}}                    % input function space
\newcommand{\U}{\mathcal{U}}                    % output function space
\newcommand{\D}{D}                              % spatial domain
\newcommand{\F}{\mathcal{F}}                    % Fourier transform
\let\Finv\relax                                % (amssymb defines \Finv)
\newcommand{\Finv}{\mathcal{F}^{-1}}
\newcommand{\kmax}{k_{\max}}
\newcommand{\dv}{d_v}                           % channel width

% ---- Rhizome operator -----------------------------------------------
\newcommand{\h}{\bm{h}}                         % hidden state
\newcommand{\cond}{\bm{c}}                      % conditioning
\newcommand{\msg}{\bm{m}}                       % message
\newcommand{\Agate}{A_\theta}                   % receiver gate
\newcommand{\Bgate}{B_\theta}                   % source gate
\newcommand{\Vval}{V_\theta}                    % value map
\newcommand{\kernel}{\kappa_\theta}             % spatial kernel
\newcommand{\Niter}{T}                          % number of iterations

% ---- 21 cm physics --------------------------------------------------
\newcommand{\xHI}{x_{\mathrm{HI}}}
\newcommand{\Tb}{\delta T_{\mathrm{b}}}
\newcommand{\cmfast}{\texttt{21cmFAST}}


% ---- Draft switch: set to \draftfalse to hide all \TODO{} notes -----
\drafttrue

\begin{document}

% ---- Title block (full width, journal style) ------------------------
\twocolumn[{%
  \begin{flushleft}
    {\small\itshape Technical report \textperiodcentered{} working draft, \today}\par
    \vspace{2pt}\rule{\textwidth}{0.5pt}\par
    \vspace{18pt}
    {\LARGE\bfseries Rhizome Neural Operator}\par
    \vspace{6pt}
    {\large\itshape Technical Report: a state-dependent integral operator for 21\,cm reionization fields}\par
    \vspace{12pt}
    Haris Kapetanovic\par
    \vspace{14pt}
  \end{flushleft}
  \noindent\hfill%
  \begin{minipage}{0.88\textwidth}
    \small
    {\centering\textbf{Abstract}\par}\vspace{3pt}
    We study a recurrent neural integral operator for predicting neutral-fraction and 21\,cm brightness-temperature fields from \texttt{21cmFAST} lightcones. A learned convolution is multiplied by channel-wise source and receiver gates computed from the current hidden field at each update. This factorization makes the effective interaction depend on the field while retaining FFT evaluation and $\mathcal{O}(N\log N)$ spatial cost at fixed width and mode count. We describe its continuum form, discrete evaluation and restrictions, including the distinction between bounded hidden states and convergence. In the reported single-seed experiments, the 3-D rhizome reaches a neutral-fraction test RMSE of 0.059 against 0.068 for the existing FNO-based pipeline; with joint prediction, the errors are 0.050 and 5.8\,mK against 0.056 and 6.5\,mK. A slice-wise 2-D model with $\pm32$ cells of LOS context improves morphology further. The best multi-field model, trained on a larger set of cones, reduces bubble-size error from 1.21 to 0.56\,Mpc and front width from 6.17 to 3.89\,Mpc, and predicts a full cone in 5.5\,s on an A100. Adding a brightness-temperature power loss suppresses excess small-scale power before reionization at a 4--8\% cost in per-voxel RMSE. These results support the complete prediction pipeline; the separate contribution of evolving-state gates remains to be established by controlled ablations.\par
    \vspace{8pt}
    \textbf{Keywords:} neural operators, Fourier neural operator, operator learning, 21\,cm cosmology, epoch of reionization, emulation
  \end{minipage}\hfill\null\par
  \vspace{16pt}
}]
\thispagestyle{plain}

{\small\tableofcontents}
\vspace{6pt}

% =====================================================================
\section{Neural operators and Fourier layers}
\label{sec:neural-operators}
% =====================================================================

\subsection{Operator learning: problem setup}
\label{sec:setup}

A neural operator learns a map between fields rather than between vectors of a prescribed size. Let $\A$ and $\U$ be normed spaces of functions on domains $D$ and $D'$, and let $\G:\A\to\U$ denote the target map. Given input--output pairs $(f_j,g_j)$, with $g_j=\G(f_j)$, we seek an approximation $\Gtheta$ whose parameters can be used on different discretizations of those fields \citep{kovachki2023neural}. The domains need not be identical, although the architectures considered here predict on the input grid.

In this work, the input contains the matter-density field and, for the multi-field models, the LOS velocity. Redshift and simulation parameters specify which map is being emulated; they are conditioning variables, distinct from the learned weights $\theta$. The outputs are the neutral hydrogen fraction $\xHI$ and, where included, the brightness temperature $T_b$. The target is the map defined by the simulation and its preprocessing, rather than an exact solution operator of a specified PDE.

The function-space viewpoint separates an architecture from its numerical evaluation. It does not by itself guarantee continuity, resolution transfer or generalization. For example, $L^2$ is useful for field-error norms, whereas bounded measurable fields in $L^\infty$ are convenient for the gated interaction below. Point samples require a sufficiently regular representative; simulation cells can instead be understood as cell averages. Changing the grid must preserve the domain, physical units and meaning of every input channel.

\subsection{Neural operators as stacked kernel integrals}
\label{sec:neural-operator-layers}

A linear integral operator maps a field $f$ to a field $g$ through a kernel:
\begin{equation}
  \label{eq:kernel-integral}
  g(r)=\int_D\kappa(r,r';\theta)f(r')\dd r'.
\end{equation}
For vector-valued fields, $\kappa$ is a matrix that mixes channels. A Green's function gives one physical example of such a kernel, when the inverse of a linear differential operator admits an integral representation. A learned kernel need not be a Green's function or satisfy a differential equation.

An integral alone is linear when the kernel is fixed. A neural operator obtains a nonlinear map by combining integral interactions with pointwise transformations. A typical hidden layer is
\begin{equation}
  \label{eq:neural-operator-layer}
  v_{\ell+1}(r)=\sigma\!\left(W_\ell v_\ell(r)+b_\ell(r)
           +(\Kop_\ell v_\ell)(r)\right),
\end{equation}
with lifting $v_0=P(f)$ and projection $\Gtheta(f)=Q(v_L)$. Here $W_\ell$ mixes channels at one position, while $\Kop_\ell$ exchanges information between positions. This distinction matters: a fixed kernel does not make the complete network linear or its response independent of the input.

\subsection{The Fourier neural operator}
\label{sec:fno}

On a periodic domain, a translation-invariant kernel depends only on separation, $\kappa(r,r';\theta)=\kappa(r-r';\theta)$. Its integral is a convolution:
\begin{equation}
  \label{eq:kernel-integral-fourier}
  \Kop f=\Finv\!\left[\hat\kappa\,\F[f]\right].
\end{equation}
The Fourier neural operator (FNO) parameterizes the matrix-valued multiplier $\hat\kappa$ directly \citep{li2021fourier}. A lifting layer maps the input into $\dv$ channels, several layers of the form \cref{eq:neural-operator-layer} combine spectral and pointwise paths, and a projection produces the output (\cref{fig:fno-architecture}).

For a uniform grid, the integral is approximated by
\begin{equation}
  \label{eq:kernel-integral-discrete}
  g(r_i)=\sum_{j=1}^{N}\kappa(r_i-r_j;\theta)f(r_j)w_j,
\end{equation}
where $w_j$ is a cell-volume weight, not a one-dimensional spacing unless $d=1$. On the unit-volume domain used below, $w_j=1/N$. The FFT evaluates this discrete convolution without forming all $N^2$ pairs.

\begin{figure*}[!tp]
  \centering
  \figplaceholder[0.85\linewidth]{fno_architecture.pdf}{FNO block diagram: lift $P$ $\to$ Fourier layers ($\F$, $R_\phi$, $\Finv$ + local $W$) $\to$ projection $Q$.}
  \caption{Fourier neural operator architecture. The spectral path is linear; the pointwise activation makes the complete layer nonlinear.}
  \label{fig:fno-architecture}
\end{figure*}

\subsection{Discretization invariance and its limits}
\label{sec:discretisation}

Each spectral path retains a finite set of Fourier indices $\Lambda$:
\begin{equation}
  \label{eq:mode-truncation}
  \widehat{\Kop f}(k)=
  \begin{cases}
    R_\phi(k)\hat f(k), & k\in\Lambda,\\
    0, & k\notin\Lambda,
  \end{cases}
\end{equation}
with $R_\phi(k)\in\C^{\dv\times\dv}$. We use $\kmax$ for the mode-count setting and $\Lambda$ for the actual retained set. The implementation uses rectangular, possibly anisotropic bands rather than a spherical cutoff. On a box of side $L$, index $k$ corresponds to physical angular wavenumber $2\pi k/L$; a mode count is therefore not itself a wavenumber in $\mathrm{Mpc}^{-1}$.

The number of spectral parameters depends on $|\Lambda|$ and $\dv$, rather than on the grid size. Lifting, projection and pointwise maps also share their weights across cells. The weights can therefore be reused on finer grids of the same domain that resolve the retained modes. This is the architectural basis of the FNO's demonstrated zero-shot super-resolution \citep{li2021fourier}. Several qualifications matter here.

\paragraph{The spectral path is band-limited; the output need not be.} The pointwise path can retain high-frequency information already present in the hidden field, and nonlinear activations can generate additional frequencies. It is therefore incorrect to describe the complete FNO output as band-limited. A finer grid provides additional evaluation points, but reliable predictions on newly resolved physical scales still depend on what the model learned and on the information in its inputs.

\paragraph{Aliasing.} Nonlinear activations and products can create frequencies above the grid's Nyquist limit. Sampling folds these frequencies into the represented band, with an error that depends on the grid. Sharing parameters across resolutions is weaker than requiring the discrete evaluations to represent exactly the same continuous operator. This distinction is formalized by \citet{bartolucci2023representation}.

\paragraph{Resolution-dependent targets.} Changing the simulation resolution may change smoothing, sub-grid prescriptions and the resolved ionization fronts. Transferring between discretizations of one field and predicting a new simulation at finer resolution are different tests. An architecture that supports the former need not succeed at the latter.

\paragraph{Boundary conditions.} The FFT evaluates a periodic convolution. This matches coeval periodic boxes and the transverse directions of the lightcones. Along the non-periodic LOS, padding and cropping define a different bounded-domain operation. Padding reduces wrap-around effects only over the separations it actually separates; it does not impose a physical boundary condition by itself.

\subsection{Limitations of a stationary spectral interaction}
\label{sec:fno-limitations}

The FNO offers an efficient global interaction. For ionization fields, the questions are whether its retained modes and pointwise nonlinearities can represent sharp boundaries and whether the learned response adapts sufficiently to heterogeneous regions.

\paragraph{Stationarity of the linear kernel.} A spectral path applies the same linear kernel at every position. The full network can still have an input-dependent response: nonlinearities, conditioning and repeated layers change its Jacobian. What a standard spectral path lacks is an explicit source--receiver modulation of its kernel within that interaction. This is the distinction addressed by the rhizome.

\paragraph{Sharp features.} A discontinuity has a broad Fourier spectrum; for example, the coefficients of a one-dimensional step decay as $1/|k|$. Truncating a linear Fourier representation can smooth the edge and produce Gibbs ringing. Ionization maps are close to binary, so front placement is sensitive to these scales. Pointwise paths and activations can recover detail, but a small spectral band does not directly transport arbitrary high-frequency structure between cells.

\paragraph{Intermediate scales.} A standard FNO has no explicit spatial hierarchy or compact local stencil. It can represent interactions at intermediate distances through its kernel, so these scales are not absent; they simply have no dedicated branch as they do in a U-Net.

\paragraph{Cost of more modes.} A dense multiplier stores $\dv^2$ complex entries per retained mode. Its parameter count grows as $\mathcal{O}(|\Lambda|\dv^2)$, with $|\Lambda|$ growing cubically with an isotropic mode count in three dimensions. Factorizing the channel matrices can reduce this cost, but adds a restriction of its own.

These observations motivate a different interaction law; they do not establish an impossibility result for FNOs. The experiments must determine whether the proposed change improves front placement, cross-correlation and morphology under comparable inputs and training conditions.

% =====================================================================
\section{Spatial structure in neural operators and the U-Net}
\label{sec:spatial-structure}
% =====================================================================

Ionized bubbles span a wide range of sizes, while their boundaries can be sharp on the grid scale. This motivates three useful architectural choices: local interactions, a hierarchy of scales, and explicit modulation by the field. Existing models address these in different ways. This section reviews the U-Net and U-FNO before introducing the particular combination used by the rhizome.

% ---------------------------------------------------------------------
\subsection{Inductive biases: locality, scale and stationarity}
\label{sec:inductive-biases}

The spatial interaction in a neural-operator layer can often be written as a kernel integral, with pointwise maps around it. Restrictions on this interaction determine some of the architecture's inductive biases. Three properties are useful for comparison.

\paragraph{Receptive field.} A convolutional neural network (CNN) uses kernels with compact support of a few grid cells, $\kappa(r-r')=0$ for $\abs{r-r'}$ larger than the stencil. A single layer is therefore strictly local; the receptive field grows only linearly with depth. The FNO is the opposite extreme: every Fourier layer couples every point to every other point.

\paragraph{Scale.} A plain CNN or FNO processes the field at one resolution. Neither has a built-in notion of coarse and fine scales; scale-dependent behaviour has to be learned implicitly. Multi-scale architectures instead process the field at several resolutions and combine the results, so that large-scale context and small-scale detail are handled by separate parts of the network.

\paragraph{Stationarity and state dependence.} Convolutional filters and standard Fourier multipliers are fixed after training. Their surrounding nonlinearities still make the complete network's response depend on the input. Attention \citep{vaswani2017attention} and graph kernel networks \citep{li2020graph} put feature dependence directly into the interaction weights. Dense pairwise evaluation costs $\mathcal{O}(N^2)$ at fixed width, but this is not a lower bound for state dependence: sparse graph interactions and factorized or linear attention reduce that cost \citep{katharopoulos2020linear}. The rhizome uses a convolution with separable endpoint gates.

\Cref{tab:inductive-biases} summarizes these properties for the architectures discussed in this report.

\begin{table*}[!tp]
  \centering
  \caption{Spatial interactions in the architectures discussed here. ``Fixed'' refers to the learned linear filters, not to the full nonlinear response. Costs hold at fixed width; pooling and boundaries can break exact translation equivariance.}
  \label{tab:inductive-biases}
  \begin{tabular}{lcccc}
    \toprule
    Architecture & Receptive field & Base filter & Scale hierarchy & Explicit gates \\
    \midrule
    CNN          & local, grows with depth   & fixed & no  & no \\
    FNO          & global spectral path     & fixed & no  & no \\
    U-Net        & local $\to$ wide via pooling & fixed & yes & no \\
    U-FNO        & global + multi-scale local & fixed & yes & no \\
    Dense attention & global                  & adaptive & no & $\mathcal{O}(N^2)$ \\
    Rhizome      & global                    & fixed + gates & no & $\mathcal{O}(N\log N)$ \\
    \bottomrule
  \end{tabular}
\end{table*}

% ---------------------------------------------------------------------
\subsection{The U-Net}
\label{sec:unet}

The U-Net \citep{ronneberger2015unet} was introduced for biomedical image segmentation, a task that, like predicting an ionization field, requires sharp, pixel-accurate boundaries informed by large-scale context. It consists of three parts (\cref{fig:unet}).

\paragraph{Contracting path.} The encoder applies a stack of small convolutions followed by a nonlinearity, then halves the resolution by pooling, and repeats. Denoting the feature map at level $s$ by $e_s$,
\begin{equation}
  \label{eq:unet-encoder}
  e_0 = C_0(a), \qquad e_{s} = C_s\big(\mathrm{Down}(e_{s-1})\big), \quad s = 1,\dots,S,
\end{equation}
where each $C_s$ is a block of $3\times3$ convolutions with nonlinearities. The number of channels is typically doubled at each level, so the network trades spatial resolution for feature richness. Because every level operates on a grid that is coarser by a factor of two, a $3\times3$ convolution at level $s$ spans $3\cdot2^{s}$ cells of the original grid. The receptive field therefore grows exponentially with the number of levels instead of linearly with depth.

\paragraph{Bottleneck.} At the coarsest level, $e_S$, the field is represented on a grid of $n/2^{S}$ points per dimension, where $n$ is the original side length with many channels. Here a few convolutions see a large fraction of the domain and capture large-scale context.

\paragraph{Expanding path and skip connections.} The decoder upsamples step by step and, at each level, concatenates the upsampled features with the encoder features of the same resolution:
\begin{equation}
  \label{eq:unet-decoder}
  d_S = e_S, \qquad d_s = C'_s\big(\big[\,\mathrm{Up}(d_{s+1}),\ e_s\,\big]\big), \quad s = S-1,\dots,0,
\end{equation}
and a final $1\times1$ convolution maps $d_0$ to the output $u$. The skip connections $e_s\to d_s$ are the key ingredient. Pooling discards the exact position of features, and upsampling alone cannot restore it. The skip connections pass the full-resolution encoder features directly to the decoder, so the network can combine \emph{what} is present (from the coarse levels) with \emph{where} it is (from the fine levels).

\begin{figure}[!htbp]
  \centering
  \figplaceholder[\linewidth]{unet_architecture.pdf}{U-Net: contracting path,
  bottleneck, expanding path with skip connections.}
  \caption{U-Net architecture with $S=3$ downsampling steps. The contracting path (left) halves the resolution $N$ and doubles the channel count $c$ at each level; the expanding path (right) reverses this. Dashed skip connections concatenate encoder features with the upsampled decoder features of the same resolution, which restores the spatial detail lost by pooling.}
  \label{fig:unet}
\end{figure}

\paragraph{Why it captures sharp structure.} The U-Net provides an explicit hierarchy of local receptive fields. Its finest levels operate on the full grid with local kernels, so edges can be represented directly in physical space instead of through a truncated set of global modes. Its intermediate levels give it receptive fields at every scale between a few cells and the full box, which matches the broad range of bubble sizes. Finally, skip connections let high-frequency detail bypass the bottleneck entirely. U-Nets are therefore strong baselines for field-to-field emulation whenever sharp, localized features matter.

\paragraph{Limitations.} The convolution kernels are defined in grid cells, so their physical support changes with resolution. A fully convolutional U-Net can often run on a different grid size, with padding where needed, but this does not make it a consistent discretization of the same physical operator. Pooling also introduces alignment dependence, so exact equivariance to arbitrary one-cell shifts is generally lost. Large-scale context passes through coarse features; whether this is sufficient depends on the task. The linear filters are shared spatially, while the nonlinear response remains input-dependent.

% ---------------------------------------------------------------------
\subsection{Hybrid architectures: U-FNO and related}
\label{sec:ufno}

The complementary strengths of the FNO (global, resolution-independent) and the U-Net (local, multi-scale) suggest combining them. The U-FNO of \citet{wen2022ufno} does this inside the Fourier layer. In addition to the spectral path $\F^{-1}(R_\phi\cdot\F v)$ and the pointwise path $W v$ of \cref{fig:fno-architecture}, each \emph{U-Fourier layer} adds a small U-Net acting on $v$:
\begin{equation}
  \label{eq:ufno-layer}
  v_{\ell+1}(x) = \sigma\Big( W v_\ell(x) + \F^{-1}\big(R_\phi\cdot\F v_\ell\big)(x) + \mathrm{UNet}_\ell[v_\ell](x) \Big),
\end{equation}
where $\mathrm{UNet}_\ell$ denotes the U-Net path. The spectral path supplies global, low-$k$ structure, while the U-Net path restores the high-$k$, localized detail that mode truncation removes. In the original application to multiphase flow, where saturation fronts are sharp, this substantially improved accuracy over the plain FNO \citep{wen2022ufno}.

The improvement comes at a cost that matters when comparing architectures. The U-Net path adds a large number of parameters, and a U-FNO is typically much larger than an FNO of the same width and depth. In our own comparison (\cref{sec:results}), the U-FNO reference model has roughly two orders of magnitude more parameters than the matched models. Better accuracy can therefore not be attributed to the architecture alone. We treat that implementation as a dense reference baseline. Parameter matching alone is not enough to isolate architecture: inputs, losses, optimization and compute also need to be controlled. The convolutional path introduces cell-scale dependence into the U-FNO, even though its Fourier multipliers can be reused across grids.

Other hybrids follow the same pattern of adding a local or multi-scale branch next to a global spectral one. They differ in which branch is used and where it is inserted, but they combine spatial branches with fixed learned filters and nonlinear maps, as in \cref{eq:ufno-layer}.

% ---------------------------------------------------------------------
\subsection{Open problem: structure that adapts to the field}
\label{sec:adaptive-structure}

A useful question is whether an explicit, inexpensive modulation of the spatial kernel improves field prediction beyond what a stack of fixed filters and nonlinearities already learns. This is an architectural question, not a claim that standard nonlinear networks cannot adapt to their inputs.

Radiative transfer motivates such modulation: absorption depends on the intervening material, so a physical source--receiver response need not be stationary. The present targets, however, come from the semi-numerical prescriptions of \cmfast{}, not a resolved transport calculation. Neither attention based only on endpoint features nor the rhizome's endpoint gates automatically encode optical depth along a ray. Intermediate information can enter through repeated contextual updates, but a physical transport interpretation requires separate evidence.

An operator suited to this problem should combine four properties: (i) global reach, so that interactions across the whole box are possible; (ii) state-dependent coupling, so that the effective connections change as the field evolves; (iii) a cost close to $\mathcal{O}(N\log N)$, so that three-dimensional grids remain tractable; and (iv) a parameterization that is independent of the grid, so that the discretization invariance of the FNO is kept. The rhizome operator, introduced in \cref{sec:rhizome}, is designed around these four requirements.

% =====================================================================
\section{Rhizome Neural Operator approach}
\label{sec:rhizome}
% =====================================================================

The name is borrowed from botany. A rhizome is a root system without a trunk or a hierarchy: any point can sprout a connection to any other, and which connections grow depends on local conditions rather than on a fixed plan. The rhizome neural operator applies this idea to operator learning. Every grid cell can exchange information with every other cell, as in a Fourier layer, but how strongly a cell sends and how strongly it listens is decided anew at every step from the current state of the field. The connections themselves therefore change as the prediction evolves, even though the parameters are fixed after training.

This section derives the operator from the kernel integrals of \cref{sec:neural-operators}, then describes the architecture (\cref{fig:rhizome-architecture}), how it is evaluated efficiently, and the variants used in the experiments of \cref{sec:results}.

% ---------------------------------------------------------------------
\subsection{Design principles}
\label{sec:rhizome-principles}

The design follows the four requirements formulated at the end of \cref{sec:adaptive-structure}:
\begin{enumerate}
  \item \textbf{Global reach.} Interactions between cells are carried by a learned kernel $\kernel(r-r')$ on the whole periodic domain, evaluated in Fourier space as in the FNO.
  \item \textbf{State-dependent coupling.} Two learned \emph{gates} modulate the kernel at the sending and the receiving end. They are recomputed from the current hidden state at every update, so the effective coupling between two cells depends on the field at both of them.
  \item \textbf{$\mathcal{O}(N\log N)$ cost.} The gates act channel-wise on each cell, so they factor out of the integral. The interaction is then a pointwise multiplication, one convolution evaluated with the FFT, and another pointwise multiplication: the same cost as a Fourier layer.
  \item \textbf{Grid-independent parameters.} The model has no convolution stencils, no pooling and no parameters tied to particular cells. The kernel is parameterized by its Fourier coefficients on the normalized domain, and everything else acts pointwise.
\end{enumerate}
Two further choices complete the design. First, there is exactly \emph{one} mechanism for spatial communication: the gated integral. There is no separate local or convolutional branch next to it, as in the U-FNO; lifting, the update proposal and decoding all act on single cells. Second, the same update block is applied $\Niter$ times with shared weights. The model is a recurrent dynamical system on the hidden state rather than a stack of distinct layers.

% ---------------------------------------------------------------------
\subsection{Mathematical background: from kernel integrals to the rhizome}
\label{sec:rhizome-math}

The rhizome is best understood as a specific member of the family of kernel integral operators introduced in \cref{sec:neural-operator-layers}. This subsection derives it from the linear integral operator \cref{eq:kernel-integral} in three steps. First, the kernel is allowed to depend on the field it acts on. Second, that dependence is restricted to a separable form that keeps the FFT evaluation of the FNO. Third, the resulting operator is iterated. Throughout, $D$ is the periodic domain normalized to unit volume, and $v:D\to\R^{\dv}$ is a hidden field with $\dv$ channels.

\paragraph{A hierarchy of kernels.} The most general kernel integral layer acting on $v$ lets the kernel depend on position and on the field at both ends of each connection:
\begin{equation}
  \label{eq:general-kernel}
  (\Kop[v]\,v)(r) = \int_D \kappa\big(r,r',v(r),v(r');\theta\big)\,v(r')\dd r'.
\end{equation}
Different neural operators correspond to different restrictions of $\kappa$ (\cref{tab:kernel-hierarchy}). If $\kappa$ depends on $r$ and $r'$ only, \cref{eq:general-kernel} is the linear operator \cref{eq:kernel-integral}: the Green's function picture of \cref{sec:neural-operator-layers}. If it depends only on $r-r'$, it is a convolution and becomes the Fourier layer \cref{eq:kernel-integral-fourier}. If it depends on the field values, the operator becomes nonlinear in $v$, since $v$ now enters both as the integrand and through the kernel. Self-attention \citep{vaswani2017attention} is the best-known example: there $\kappa\propto\exp\!\big(q(v(r))\cdot k(v(r'))\big)$, normalized over $r'$, and evaluating it on a grid costs $\mathcal{O}(N^2)$.

The rhizome sits between the Fourier layer and attention. It keeps a stationary kernel but multiplies it by two factors, one that depends on the field at the receiving point and one that depends on the field at the sending point:
\begin{equation}
  \label{eq:rhizome-effective-kernel}
  \kappa^{\mathrm{eff}}(r,r') = \diag\!\big(A(v(r))\big)\,\kernel(r-r')\,\diag\!\big(B(v(r'))\big),
\end{equation}
with $A,B: \R^{\dv}\to(0,2)^{\dv}$ the receiver and source gates. Setting $A=B=1$ recovers the Fourier layer. Compared with attention, the dependence on the two endpoints is \emph{separable}: it is a product of a function of $v(r)$ and a function of $v(r')$. This is also the idea behind linear attention \citep{katharopoulos2020linear}, which replaces the softmax by a product of feature maps. The rhizome applies it channel by channel and multiplies it by a learned relative-position kernel $\kernel(r-r')$.

\begin{table*}[!tp]
  \centering
  \caption{Kernel integral operators ordered by what the kernel may depend on. $N$: number of grid points.}
  \label{tab:kernel-hierarchy}
  \small
  \setlength{\tabcolsep}{4pt}
  \begin{tabular}{@{}llccc@{}}
    \toprule
    Operator & Kernel $\kappa$ & Linear & Stationary & Cost \\
    \midrule
    linear (Green's function) & $\kappa(r,r')$ & yes & no & $\mathcal{O}(N^2)$ \\
    Fourier layer & $\kappa(r-r')$ & yes & yes & $\mathcal{O}(N\log N)$ \\
    \textbf{rhizome} & $\diag A(v(r))\,\kernel(r-r')\,\diag B(v(r'))$ & no & no & $\mathcal{O}(N\log N)$ \\
    attention & $\kappa(v(r),v(r'))$, normalized & no & no & $\mathcal{O}(N^2)$ \\
    \bottomrule
  \end{tabular}
\end{table*}

\paragraph{Operator form.} Write $M_a$ for the multiplication operator $(M_a f)(r)=a(r)\,f(r)$, applied channel by channel, and $C_\kappa$ for the convolution $(C_\kappa f)(r)=\int_D\kappa(r-r')f(r')\dd r'$. The rhizome message then reads
\begin{equation}
  \label{eq:rhizome-operator-form}
  \msg = M_{A(v)}\;C_{\kernel}\;M_{B(v)}\;V(v),
\end{equation}
where $V$ is the pointwise value map. Each factor is simple in one basis: the multiplication operators are diagonal in position space, and the convolution is diagonal in Fourier space, where $\widehat{C_\kappa f}(k)=\hat\kappa(k)\hat f(k)$. The rhizome therefore alternates between the two bases: multiply in position space, transform, multiply in Fourier space, transform back, multiply again. This is the same structure as the split-step Fourier methods of computational physics, where a potential acts in position space and the kinetic term in momentum space, each where it is diagonal. Like them, it costs two FFTs per application.

The factorization also makes the operator well behaved. Since the gates lie in $(0,2)$, both multiplication operators have norm at most 2 on $L^2(D)$. By Parseval's theorem, the norm of the convolution is $\sup_k\norm{\hat\kernel(k)}$. Hence
\begin{equation}
  \label{eq:rhizome-bound}
  \norm{M_{A}\,C_{\kernel}\,M_{B}\,f}_{L^2} \;\le\; 4\,\sup_{\abs{k}\le\kmax}\norm{\hat\kernel(k)}\;\norm{f}_{L^2},
\end{equation}
independently of the field that produced the gates. State dependence cannot make the interaction blow up.

\paragraph{Mode coupling: smooth transport, sharp modulation.} \Cref{sec:discretisation} noted that the spectral path of an FNO is band-limited: everything above $\kmax$ in its output has to be generated by the pointwise nonlinearity. In the rhizome, the gates change this. A product in position space is a convolution in Fourier space, so the Fourier transform of \cref{eq:rhizome-operator-form} is
\begin{equation}
  \label{eq:rhizome-mode-coupling}
  \begin{aligned}
  \hat\msg(k) &= \sum_{q}\hat A(k-q)\;\hat\kernel(q)\;\widehat{(B\,V)}(q),\\
  \hat\kernel(q) &= 0 \quad\text{for } \abs{q}>\kmax .
  \end{aligned}
\end{equation}
The kernel only transports the band $\abs{q}\le\kmax$. Even so, the source gate first folds all input wavenumbers into that band, through $\widehat{BV}(q)=\sum_p\hat B(q-p)\hat V(p)$, and the receiver gate then spreads the transported band over every wavenumber present in $A$. If the receiver gate is sharp, for example close to zero inside an ionized region and close to two outside, a smooth message is multiplied by a sharp mask, and the result has a sharp edge exactly where the gate changes. The rhizome can thus carry information smoothly over long distances with few modes while placing sharp structure where the state of the field puts it. Whether the trained gates actually behave like this has not been checked directly; \cref{eq:rhizome-mode-coupling} shows only that the mechanism is available.

\paragraph{Iteration as a path series.} A single application of \cref{eq:rhizome-effective-kernel} couples $r'$ to $r$ through the gates at the two endpoints only. Iteration adds paths. Ignoring for a moment the nonlinearity and the damping of the update, and freezing the gates, the recurrence becomes $v^{t+1}=c+\Kop v^t$ with $\Kop=M_A C_{\kernel} M_B$. After $\Niter$ steps this gives
\begin{equation}
  \label{eq:neumann}
  v^{\Niter} = \sum_{n=0}^{\Niter-1}\Kop^{\,n}\,c \;+\; \Kop^{\Niter} v^0,
\end{equation}
the truncated Neumann series of the solution $v=(I-\Kop)^{-1}c$ of the Fredholm integral equation of the second kind $v=c+\Kop v$. The resolvent $(I-\Kop)^{-1}$ plays the role of the Green's function of \cref{sec:neural-operator-layers}, now for a state-dependent problem. Each term of the series is a sum over paths: the second-order term, for example,
\begin{multline}
  \label{eq:two-hop}
  (\Kop^2 f)(r) = \int_D\!\int_D A(r)\,\kernel(r-r_1)\,\underbrace{B(r_1)A(r_1)}_{\text{gates at } r_1}\\
  \times\kernel(r_1-r_2)\,B(r_2)\,f(r_2)\dd r_1\dd r_2,
\end{multline}
couples $r_2$ to $r$ through every intermediate point $r_1$, weighted by the gates there. This is the mathematical form of the path effects mentioned in \cref{sec:rhizome-interaction}. A region whose gates are small blocks the paths through it and so shadows what lies behind it, even though no single update can see beyond its two endpoints. In scattering theory the same expansion is the Born series, with each term adding one more scattering event. The actual rhizome differs from this linear picture through the $\tanh$ proposal, the damping $(1-\Delta t)$ and the gates being recomputed at every step, but the path structure is the same.

\paragraph{Continuous-time view.} The update of \cref{alg:rhizome}, $\h^{t+1}=(1-\Delta t)\h^t+\Delta t\,\bm p^t$, is the forward Euler discretization of the nonlocal evolution equation
\begin{equation}
  \label{eq:rhizome-ode}
  \partial_\tau \h(r,\tau) = -\h(r,\tau) + \tanh\!\Big(\Phi\big(\h(r,\tau),\ (\Kop[\h]\,V(\h))(r),\ \cond(r)\big)\Big),
\end{equation}
integrated from $\tau=0$ to $\tau=\Niter\Delta t$ (3 in the experiments), with $\Phi$ the pointwise proposal network. Equations of this form, a decay term plus a saturating function of a spatially integrated input, are the neural field equations of computational neuroscience \citep{amari1977dynamics}, in which $\partial_t u=-u+\int w(r-r')\,S(u(r'))\dd r'+I(r)$ describes the activity of a continuous sheet of neurons. The rhizome generalizes that picture with many channels, a learned kernel and connections gated by the state. Viewed this way, it is also a neural ODE \citep{chen2018neural} on a function space, integrated with a fixed explicit scheme. The bound on the hidden state follows directly: if $\abs{\h}\le1$ componentwise, then so is $(1-\Delta t)\h+\Delta t\,\bm p$ for $\abs{\bm p}\le1$ and $0<\Delta t\le1$.

\paragraph{Discretization.} On a uniform grid of $N$ points on the unit-volume domain, the integral in \cref{eq:general-kernel} becomes the sum \cref{eq:kernel-integral-discrete} with weights $\Delta r_j=1/N$. With the kernel represented by its retained Fourier modes, $\kernel(r)=\sum_{\abs{k}\le\kmax}\hat\kernel(k)\,e^{2\pi i k\cdot r}$, this sum is evaluated \emph{exactly} by the FFT:
\begin{equation}
  \label{eq:fft-quadrature}
  \begin{aligned}
  \frac1N\sum_{j=1}^{N}\kernel(r_i-r_j)\,g(r_j) &= \sum_{\abs{k}\le\kmax}\hat\kernel(k)\,\hat g_k\,e^{2\pi i k\cdot r_i},\\
  \hat g_k &= \frac1N\sum_{j=1}^N g(r_j)\,e^{-2\pi i k\cdot r_j},
  \end{aligned}
\end{equation}
provided every dimension of the grid has enough points (at least $2\kmax$ in the implementation) for the retained modes to remain distinct. The only approximation relative to the continuum operator is the quadrature itself. This identity is what the numerical verification in \cref{sec:rhizome-discretisation} checks, by evaluating the left-hand side directly over all pairs.

% ---------------------------------------------------------------------
\subsection{State-dependent integral interaction}
\label{sec:rhizome-interaction}

Each grid cell $i$ at position $r_i$ carries a hidden state $\h_i^t$ of $\dv$ channels and a fixed conditioning vector $\cond_i$, built from the input density and the global parameters. At update $t$, cell $i$ receives the message
\begin{equation}
  \label{eq:rhizome-message}
  \begin{split}
  \msg_i^t = \frac{1}{N}\sum_{j=1}^{N}
  &\diag\!\big(\Agate(\h_i^t,\cond_i)\big)\,
  \kernel(r_i-r_j)\\
  &\times\diag\!\big(\Bgate(\h_j^t,\cond_j)\big)\,
  \Vval(\h_j^t,\cond_j),
  \end{split}
\end{equation}
where
\begin{itemize}
  \item $\Vval$ builds the \emph{value} that cell $j$ transmits from its state and conditioning,
  \item $\Bgate$ is the \emph{source gate}: how strongly cell $j$ transmits, per channel,
  \item $\kernel(r_i-r_j)\in\R^{\dv\times\dv}$ is a learned, periodic, channel-mixing kernel that carries the signal across the separation $r_i-r_j$,
  \item $\Agate$ is the \emph{receiver gate}: how strongly cell $i$ accepts what arrives, per channel,
  \item $N$ is the number of grid cells, so that $1/N$ is the uniform quadrature weight on the normalized periodic box.
\end{itemize}

\begin{figure*}[!tp]
  \centering
  \figplaceholder[0.92\linewidth]{rhizome_architecture.pdf}{Rhizome block diagram.}
  \caption{Rhizome Neural Operator architecture. \emph{Top:} the input density bands and the global scalars (redshift $z$, simulation parameters $\vartheta$) are lifted and embedded into a conditioning field $\cond$, which initializes the hidden state and is fed into every update. One update block with shared weights is applied $\Niter$ times; a pointwise decoder maps the final state to $\xHI$ (and, in the multi-field model, to $T_b$). \emph{Bottom:} one update. The gates produce a source factor $\Bgate$ and a receiver factor $\Agate$ from the current state. The source-gated value is convolved with the kernel $\kernel$ in Fourier space and multiplied by the receiver gate, giving the message $\msg^t$. A pointwise proposal combines message, state and conditioning, and an explicit Euler step updates the state.}
  \label{fig:rhizome-architecture}
\end{figure*}

Read from right to left, \cref{eq:rhizome-message} describes a transmission: cell $j$ prepares a value, decides how loudly to send it, the kernel carries it over the distance $r_i-r_j$, and cell $i$ decides how much of it to take in. It is the quadrature of the operator \cref{eq:rhizome-operator-form}, with the effective kernel \cref{eq:rhizome-effective-kernel} evaluated at the current state $(\h^t,\cond)$. This is the key difference from the FNO: the base kernel $\kernel$ is stationary, but the effective kernel depends on the field at both ends of every connection and changes from update to update.

The gates are deliberately simple. They act channel by channel on single cells and therefore factor out of the sum in \cref{eq:rhizome-message}. This is what keeps the cost at that of a Fourier layer: the network never forms the $N\times N$ matrix of pairwise couplings, which for a $140\times140$ slice alone would have about $3.8\times10^8$ entries per channel pair. The price is expressivity. The coupling between two cells is a product of a sender factor and a receiver factor around a fixed kernel, not an arbitrary function of both states as in attention. In particular, a single update cannot make the coupling between $r$ and $r'$ depend on the cells \emph{between} them. Such path effects, like the shadowing of a region by neutral gas along the line of sight, emerge only over several updates, through the path series \cref{eq:two-hop}.

% ---------------------------------------------------------------------
\subsection{Gating: source and receiver factors}
\label{sec:rhizome-gates}

A single pointwise network maps $[\h^t,\cond]$ to $2\dv$ logits per cell, which are split into the source and receiver gates and passed through $2\,\mathrm{sigmoid}(\cdot)$. Each gate value therefore lies in $(0,2)$ and equals one at zero logit: an untrained gate leaves the kernel unchanged, and training can then learn to amplify or suppress transmission and reception.

Three properties of the gates are worth stating explicitly:
\begin{itemize}
  \item \textbf{They are recomputed at every update} from the previous state and the fixed conditioning. A cell that becomes ionized during the rollout can change how it sends and receives in the next update.
  \item \textbf{Source and receiver are independent.} Interactions may be asymmetric: cell $j$ can influence cell $i$ strongly while $i$ barely influences $j$. Since the kernel is not constrained to be positive, interactions may also be signed.
  \item \textbf{They are channel-wise.} Each of the $\dv$ channels has its own gate, so different channels can carry different kinds of information with different reach. Channels are mixed only by the kernel $\kernel$ and by the pointwise networks.
\end{itemize}
No physical constraint is imposed: the operator does not enforce reciprocal couplings, conservation of photons or any form of radiative transfer. Whatever physical interpretation the gates acquire is learned from the data.

% ---------------------------------------------------------------------
\subsection{Recurrent time-stepping}
\label{sec:rhizome-recurrence}

The hidden state is initialized from the conditioning and updated $\Niter$ times with the same weights (\cref{alg:rhizome}). Each update forms the message, combines it with the state and the conditioning in a pointwise \emph{proposal}, and takes an explicit Euler step of size $\Delta t$ towards that proposal. The final state is decoded pointwise into the $\xHI$ logit.

\begin{algorithm}[!htb]
  \caption{Rhizome operator forward pass}
  \label{alg:rhizome}
  \begin{algorithmic}[1]
    \Require density bands $\delta$, redshift $z$, parameters $\vartheta$
    \State $\cond \gets \mathrm{lift}(\delta) + \mathrm{embed}(z,\vartheta)$
    \State $\h^0 \gets \tanh(\cond)$
    \For{$t = 0,\dots,\Niter-1$} \Comment{shared weights}
      \State $(\bm{s},\bm{r}) \gets 2\,\mathrm{sigmoid}\big(\mathrm{gates}([\h^t,\cond])\big)$ \Comment{source, receiver}
      \State $\bm{v} \gets \Vval([\h^t,\cond])$
      \State $\msg^t \gets \bm{r}\odot \Finv\!\big(R\cdot\F(\bm{s}\odot\bm{v})\big)$
      \State $\bm{p} \gets \tanh\big(\mathrm{pointwise}([\h^t,\msg^t,\cond])\big)$
      \State $\h^{t+1} \gets (1-\Delta t)\,\h^t + \Delta t\,\bm{p}$
    \EndFor
    \State \Return $\xHI = \mathrm{sigmoid}\big(\mathrm{decode}(\h^{\Niter})\big)$
  \end{algorithmic}
\end{algorithm}

Because the proposal is bounded by $\tanh$ and the update is a convex combination for $0<\Delta t\le1$, the hidden state stays bounded by one throughout the rollout. This keeps long rollouts numerically stable, but it is not a convergence guarantee: nothing forces the iteration towards a fixed point. In practice the model is trained and evaluated at a fixed number of updates. All experiments in \cref{sec:results} use $\Niter=6$ and $\Delta t=0.5$.

Weight sharing makes the parameter count independent of $\Niter$. It also makes the model behave like a learned dynamical system rather than a deep feed-forward network, which raises the question of what happens away from the training horizon. A small pilot on synthetic $32\times32$ lightcones answers this (\cref{tab:rhizome-pilot}). The tied rhizome improved modestly on an earlier additive recurrent model (Fourier message plus separate local message) at an equal parameter count, while an untied version with independent weights per update, at four times the parameters, was more accurate still. Running the tied model for more updates than it was trained with degraded its accuracy, although the hidden state stayed bounded. The number of updates is therefore part of the trained model, not a free accuracy knob.

\begin{table}[!htbp]
  \centering
  \caption{Architecture pilot on synthetic $32\times32$ lightcones (single seed, CPU; 512 training slices, 64 test slices). \emph{Top:} rhizome against the additive recurrent baseline. \emph{Bottom:} the tied rhizome evaluated at different numbers of updates; it was trained with four. Mixed slices have mean $\xHI$ between 0.05 and 0.95. This toy problem is not a physical simulation; the table characterizes the architecture, not its accuracy on 21\,cm data.}
  \label{tab:rhizome-pilot}
  \begin{tabular}{lrr}
    \toprule
    Model & Params & RMSE, mixed \\
    \midrule
    Rhizome, tied           & 15.6\,k & 0.196 \\
    Additive recurrent      & 15.9\,k & 0.207 \\
    Rhizome, untied         & 60.3\,k & 0.181 \\
    \bottomrule
  \end{tabular}
  \\[8pt]
  \begin{tabular}{rr}
    \toprule
    Updates $\Niter$ & RMSE, mixed \\
    \midrule
    1 & 0.260 \\
    2 & 0.205 \\
    \textbf{4} (trained) & \textbf{0.196} \\
    8 & 0.265 \\
    \bottomrule
  \end{tabular}
\end{table}

% ---------------------------------------------------------------------
\subsection{Discretization and spectral evaluation}
\label{sec:rhizome-discretisation}

The FFT in \cref{fig:rhizome-architecture} is not a separate component that predicts anything; it is simply an efficient and exact way to evaluate the sum over $j$ in \cref{eq:rhizome-message} (\cref{eq:fft-quadrature}). The kernel $\kernel$ is parameterized by its Fourier coefficients $R(k)$ for $\abs{k}\le\kmax$, exactly as in the FNO (\cref{eq:mode-truncation}). Three details make this evaluation correct:
\begin{itemize}
  \item \textbf{Quadrature.} With an orthonormal FFT pair, the forward and inverse transforms together already produce the $1/N$ weight of \cref{eq:rhizome-message}. Adding a further factor for the cell area would make messages shrink as the resolution grows.
  \item \textbf{Real kernel.} The zero mode and the conjugate-symmetric modes are constrained so that $\kernel$ is real in physical space.
  \item \textbf{Verification.} The FFT implementation was checked against an independent calculation of the left-hand side of \cref{eq:fft-quadrature} that sums the retained Fourier series over all source--receiver pairs explicitly, without any FFT. Outputs agree to about $10^{-15}$, and so do the gradients with respect to inputs and parameters.
\end{itemize}

The architecture inherits the FNO's discretization properties (\cref{sec:discretisation}). Since no parameter refers to a particular cell, shifting the input periodically shifts the output in the same way. The same weights can be applied to any grid that represents the same physical box with at least $2\kmax$ points per dimension. The caveats also carry over: the gates and the proposal are nonlinear and alias on coarse grids, so transfer across resolutions has to be validated rather than assumed. It has not been tested in this work.

% ---------------------------------------------------------------------
\subsection{Model variants}
\label{sec:rhizome-variants}

The same operator was used in four configurations, summarized in \cref{tab:rhizome-variants}. They differ in what the model sees along the line of sight (LOS) and in which fields it predicts.

\begin{table*}[!tp]
  \centering
  \caption{Rhizome variants used in the experiments. LOS: line of sight.}
  \label{tab:rhizome-variants}
  \small
  \begin{tabular}{@{}p{2.6cm}p{4.4cm}p{3.6cm}p{3.6cm}@{}}
    \toprule
    Variant & Input & Output & Kernel \\
    \midrule
    2-D slice
      & one transverse slice + 13 LOS density bands ($\pm32$ cells)
      & $\xHI$ on the slice
      & 2-D, 24 modes \\
    3-D window
      & native LOS windows of 256 cells + $2\times32$ halo
      & $\xHI$ in the window
      & 3-D, modes (36,\,36,\,24), rank-48 factorized \\
    Slice-wise 2-D
      & 2-D slice model applied to every slice of a cone
      & full 3-D $\xHI$ cone, restacked
      & as 2-D slice \\
    Multi-field slice-wise
      & density and LOS-velocity bands, global-history band means
      & $\xHI$ and $T_b$ (structured head)
      & as 2-D slice \\
    \bottomrule
  \end{tabular}
\end{table*}

\paragraph{2-D slices with LOS context.} The first variant predicts one transverse $140\times140$ slice of a lightcone. Its input is the density on that slice and on the neighbouring slices: single slices within $\pm3$ cells and band means further out, up to $\pm32$ cells, for 13 density channels in total. The conditioning scalars are $\log(1+z)$ and the eleven simulation parameters.

\paragraph{3-D native windows.} The second variant works directly in three dimensions on windows of the native lightcone. It uses anisotropic mode counts (36 in each transverse direction, 24 along the LOS) and a rank-48 factorized spectral kernel to keep the parameter count at 5.8\,M. Because the LOS is not periodic, the windows are zero-padded by 64 cells before the FFT. A code review found that this padding is too short to separate all physical LOS offsets: displacements of $+160$ and $-160$ cells fall on the same residue of the padded circle, so the base kernel cannot treat them differently. The gates can still modulate these couplings, but the 3-D kernel is a constrained rather than a fully general LOS interaction.

\paragraph{Slice-wise reconstruction of 3-D cones.} The third variant applies the 2-D slice model to every native slice of a lightcone, each with its own $\pm32$-cell LOS context, and restacks the predictions into a full 3-D cone (\cref{fig:slicewise}). In the language of windowed inference this is maximal overlap: window one slice, stride one slice. The model never sees the full LOS jointly; it only needs to know what lies within about 46\,Mpc of each slice. As \cref{sec:results} shows, this turns out to be enough, and it avoids both the padding issue above and the memory cost of 3-D windows.

\begin{figure}[!htbp]
  \centering
  \figplaceholder[\linewidth]{slicewise_reconstruction.pdf}{Slice-wise reconstruction.}
  \caption{Slice-wise reconstruction of a 3-D lightcone. For each slice $i$ along the line of sight, the 2-D rhizome receives 13 density channels: the seven single slices within $\pm3$ cells and six band means out to $\pm32$ cells (blue). In the multi-field model the same bands are built for the LOS velocity. Global scalars enter as conditioning. The prediction for slice $i$ is computed independently, and the predictions for all slices are restacked into the 3-D cone.}
  \label{fig:slicewise}
\end{figure}

\paragraph{Multi-field prediction.} The final variant predicts both $\xHI$ and the 21\,cm brightness temperature $T_b$. It receives 13 LOS bands each of density and LOS velocity. Its scalars are $1/(1+z)$, the eleven parameters and the band means of two emulated global histories, $\bar x_{\mathrm{HI}}(z)$ and $\bar T_b(z)$. The model outputs an $\xHI$ logit and a spin variable $u$, and $T_b$ is assembled by a \emph{structured head} that encodes the known form of the signal:
\begin{equation}
  \label{eq:structured-head}
  T_b = T_0(z)\,\xHI\,(1+\delta)\,F_v\,\big(1-e^{u}\big),
\end{equation}
where $T_0(z)$ is the redshift-dependent amplitude, $F_v$ the correction for the LOS velocity gradient, and $1-e^{u}$ plays the role of the spin-temperature factor $1-T_{\mathrm{CMB}}/T_S$. The physical prefactor $T_0(1+\delta)F_v$ is computed exactly, so the network only has to learn the neutral fraction and the spin factor.

% ---------------------------------------------------------------------
\subsection{Assumptions and limitations}
\label{sec:rhizome-assumptions}

The rhizome makes fewer structural assumptions than the architectures of \cref{sec:spatial-structure}, but it does make some:
\begin{itemize}
  \item \textbf{Periodicity.} The kernel lives on a periodic domain. This is exact in the transverse directions of the lightcones, which are periodic, but not along the LOS. The 2-D and slice-wise variants avoid the issue by treating the LOS through input bands; the 3-D variant relies on zero padding (\cref{sec:rhizome-variants}).
  \item \textbf{Product-form coupling.} The effective kernel \cref{eq:rhizome-effective-kernel} is a stationary kernel modulated at its two ends. It cannot represent couplings that depend jointly and arbitrarily on the states of both cells, or on the cells between them, within a single update.
  \item \textbf{No physics in the operator.} Reciprocity, conservation and radiative transfer are not imposed. Physics enters only through the inputs, the conditioning and, for $T_b$, the structured head.
  \item \textbf{Fixed rollout length.} The model is accurate at the number of updates it was trained with, not beyond (\cref{tab:rhizome-pilot}). There is no convergence or fixed-point guarantee.
  \item \textbf{Untested resolution transfer.} Although the parameters are grid-independent, all results in this report were obtained at a single resolution.
\end{itemize}

% ---------------------------------------------------------------------
\subsection{Relation to FNO, U-FNO and attention}
\label{sec:rhizome-relations}

\Cref{tab:inductive-biases} places the rhizome among the architectures of \cref{sec:spatial-structure}. Compared with the \textbf{FNO}, it keeps the global spectral kernel and its $\mathcal{O}(N\log N)$ cost, adds the gates that make the effective kernel state-dependent, and replaces the stack of independent layers by a weight-shared recurrence. The pointwise proposal takes over the role of the FNO's pointwise path $W$. With gates fixed to one and untied weights, an update reduces to a Fourier layer with a nonlinear pointwise path.

Compared with the \textbf{U-Net and U-FNO}, it has no convolution stencils, no pooling and no explicit hierarchy of scales. Multi-scale behaviour has to come from the kernel, which spans all scales up to $\kmax$, and from the repeated updates. This keeps the parameterization independent of the grid, which the convolutional path of the U-FNO does not.

Compared with \textbf{attention}, it shares the central idea that the coupling between two points should depend on their content, but it restricts that dependence to a product of per-cell factors around a stationary kernel. This restriction is what brings the cost down from $\mathcal{O}(N^2)$ to $\mathcal{O}(N\log N)$: the rhizome trades the full flexibility of pairwise attention for a form that can be evaluated with a single FFT pair per update.

The cost of a forward pass is $\Niter$ times that of one Fourier layer, $\mathcal{O}(\Niter\,\dv\,N\log N)$ for the transforms plus pointwise work, and the parameter count is that of a single update block.

% =====================================================================
\section{Training and Results}
\label{sec:results}
% =====================================================================

This section reports how the rhizome was trained and how it compares with the best models of the existing \texttt{fno-21cm} pipeline on the task of predicting the neutral-fraction field $\xHI$, and later also the brightness temperature $T_b$, from the matter density of \cmfast{} lightcones. The experiments followed the order of \cref{tab:rhizome-variants}: 2-D slices, native 3-D windows, multi-field 3-D, and finally slice-wise reconstruction of full cones. \Cref{fig:metric-ratios} summarizes the outcome.

\begin{figure}[!htbp]
  \centering
  \figplaceholder[\linewidth]{metric_ratios.pdf}{Headline metrics relative to fno.}
  \caption{Headline results on 199 held-out multi-field test cones. Each bar is the error of a rhizome model divided by the error of the best \texttt{fno-21cm} multi-field model (\texttt{tbs+histtb}); values below the dashed line are better than the baseline. Front width is measured as the excess over the true front width of 2.77\,Mpc. The reported slice-wise errors are lower on the listed metrics except the small-scale $T_b$ power before reionization, which a power term in the loss reduces (\cref{sec:results-absorption}). Numbers in \cref{tab:main-results}.}
  \label{fig:metric-ratios}
\end{figure}

% ---------------------------------------------------------------------
\subsection{Datasets}
\label{sec:datasets}

All data come from one suite of 6,600 \cmfast{} \citep{mesinger2011cmfast} simulations (version 4.1.1), each with eleven sampled simulation parameters. Every simulation provides a 3-D lightcone of $140\times140$ transverse cells over a 200\,Mpc box (1.43\,Mpc per cell). The native line-of-sight length depends on cosmology and is between 2,074 and 2,919 cells at the same spacing. The lightcones contain the density, the neutral fraction, the brightness temperature and the LOS velocity, among other fields. Two properties of the data matter for the models:
\begin{itemize}
  \item Each lightcone reuses the periodic 200\,Mpc spatial realization along the LOS, with a tiling period of 140 cells. The fields can evolve with redshift along that direction, so this is repeated spatial sampling rather than an identical copy of the evolving field.
  \item Ionized cells have $\xHI$ exactly zero; partially ionized values appear mostly where the lightcone interpolates between two simulation snapshots.
\end{itemize}
33 simulations (0.5\%) contain non-finite $T_b$ values and are excluded from every multi-field split. \Cref{tab:datasets} lists the splits used.

\begin{table*}[!tp]
  \centering
  \caption{Training, validation and test splits. All test sets are held out by simulation (cone), never by slice.}
  \label{tab:datasets}
  \small
  \begin{tabular}{@{}p{2.9cm}p{4.3cm}p{2.6cm}p{4.2cm}@{}}
    \toprule
    Split & Used for & Train / val / test cones & Samples \\
    \midrule
    2-D slice cache & 2-D sweep and first 2-D model & 5,280 train (own split) & 16 window + 2 background slices per cone \\
    $\xHI$ (fno split) & 3-D rhizome; fair slice-wise model & 1,600 / 200 / 200 & 3-D windows; or 48 + 6 slices per cone \\
    Multi-field & 3-D multi-field; slice-wise multi-field & 1,594 / 199 / 199 & as above, cleaned $T_b$ \\
    Multi-field, all cones & best slice-wise multi-field model & 6,169 / 199 / 199 & 16 + 2 slices per cone (98.5\,k) \\
    \bottomrule
  \end{tabular}
\end{table*}

The $\xHI$ and multi-field splits are those of the \texttt{fno-21cm} pipeline. All comparisons with \texttt{fno} models below use identical test cones and the same evaluation code.

% ---------------------------------------------------------------------
\subsection{Training protocol}
\label{sec:training}

\paragraph{Losses.} $\xHI$ is trained with binary cross-entropy (BCE) on the logit. In the multi-field models, $T_b$ is added as a mean-squared error on the normalized field, through the structured head of \cref{eq:structured-head}. The choice of BCE turned out to matter beyond the operator (\cref{sec:results-3d-mf}). The \texttt{fno} baselines train $\xHI$ with a mean-squared error through a sigmoid, whose gradient vanishes when the prediction saturates.

\paragraph{Optimization and stopping.} The learning rate is $2\times10^{-3}$ throughout; the 3-D models use warmup and cosine decay. The 2-D models are stopped by a predefined plateau rule on the validation loss, and the checkpoint with the best validation score is kept. The 3-D models are trained for 20 epochs of 6,400 windows and then continued with a short \emph{warm restart} (10 epochs at $5\times10^{-4}$ with cosine decay), which improved test RMSE by 3--6\%, against about 1\% for a continuation at constant small learning rate. The 2-D models use the dihedral symmetries of the square slice as data augmentation.

\paragraph{Meaning of the neutral-fraction loss.} BCE accepts targets in $[0,1]$, including partially ionized simulation cells; it need not mean that the output is a Bernoulli sample. For deterministic targets it acts as a bounded-fraction regression loss. If the compressed inputs leave some morphology unresolved, both BCE and squared error are minimized pointwise by the conditional mean of the target. BCE alone therefore does not eliminate averaging of uncertain boundaries or produce a distribution of possible fields.

\paragraph{Hyperparameters.} \Cref{tab:hyperparameters} lists the main settings. Every model is a single training run with seed 0, except for one repeat of the main multi-field configuration with seed 1 (\cref{sec:results-variations}).

\paragraph{Recorded training details of the slice-wise multi-field models.} The run metadata fix the following settings. The optimizer is Adam without weight decay or gradient clipping. Validation runs every 1,000 steps. The learning rate is halved after four validations without a relative improvement of 0.2\%, down to $2\times10^{-5}$; training stops after eight such validations at the minimum rate, but not before 20,000 steps. The checkpoint with the lowest validation score, the mean of the normalized $\xHI$ and $T_b$ mean-squared errors, is kept. The loss is BCE on $\xHI$ plus the normalized $T_b$ MSE with weight 1. Density is divided by 10, $\xHI$ is not rescaled, and $T_b$ is normalized as $(T_b+28.1\,\mathrm{mK})/56.1\,\mathrm{mK}$ with training-set statistics. Per cone, 16 slices are drawn stratified in the global history within $0.02\le\bar x_{\mathrm{HI}}\le0.98$, plus two background slices outside that window, over $5.0\le z\le24.97$. The all-cones cache gives 98,509 training and 3,080 validation slices; the slice-level test set has 3,129 slices from 199 cones.

\begin{table*}[!tp]
  \centering
  \caption{Main hyperparameters.}
  \label{tab:hyperparameters}
  \small
  \begin{tabular}{@{}p{3.6cm}p{5.0cm}p{5.4cm}@{}}
    \toprule
    & 2-D slice / slice-wise models & 3-D window model \\
    \midrule
    width $\dv$ & 64 (128 in early runs) & 48 \\
    Fourier modes $\kmax$ & 24 & (36, 36, 24), rank-48 factorized \\
    updates $\Niter$ / step $\Delta t$ & 6 / 0.5 & 6 / 0.5 \\
    batch / sample & 32 slices & native window, $140\times140\times256$, including $2\times32$ halo \\
    learning rate & $2\times10^{-3}$, plateau-stopped & $2\times10^{-3}$ (warmup + cosine), 20 epochs + warm restart \\
    precision & float32 & bfloat16 autocast, float32 FFT \\
    parameters & grow as $\dv^2$: 1.1\,M (w32) to 18.3\,M (w128) & 5.8\,M \\
    training time & 2.8--3.5\,h (A100) & $\sim$12\,h + 6\,h restart (H200) \\
    \bottomrule
  \end{tabular}
\end{table*}

% ---------------------------------------------------------------------
\subsection{Evaluation metrics}
\label{sec:metrics}

Per-voxel RMSE is reported, but it is a weak guide to physical quality: a model can lower its RMSE by blurring fronts. The evaluation therefore uses the physical suite of the \texttt{fno-21cm} pipeline, computed on full test cones. Most metrics are evaluated in the \emph{active} stage of reionization, slices with mean neutral fraction between 0.05 and 0.95, and some also \emph{before reionization}, $\bar x_{\mathrm{HI}}\ge0.98$, where $T_b$ is in absorption. \Cref{tab:metrics} lists them.

\begin{table*}[!tp]
  \centering
  \caption{Evaluation metrics. All are computed per test cone and summarized by the median over cones.}
  \label{tab:metrics}
  \small
  \begin{tabular}{@{}p{4.3cm}p{9.6cm}@{}}
    \toprule
    Metric & What it measures \\
    \midrule
    $|\Delta\tau_e|/\sigma_{\mathrm{Planck}}$ & error in the CMB optical depth implied by the predicted global history, in units of the Planck uncertainty \\
    $|\Delta z|$ at $\bar x=0.25/0.5/0.75$; duration & timing of reionization: error in the redshift at which the global neutral fraction crosses each value, and in the duration \\
    $\xHI$ PDF $W_1$ & Wasserstein-1 distance between predicted and true one-point distributions of $\xHI$ \\
    partially ionized fraction & fraction of cells with intermediate $\xHI$ (truth $\approx0.09$); measures hedging \\
    $b_1$ pull; $P_\varepsilon$ ratio & EFT diagnostics: error in the fitted linear bias in units of its uncertainty; ratio of stochastic power to the truth (ideal 1) \\
    bubble-size $W_1$ & Wasserstein-1 distance (Mpc) between mean-free-path bubble-size distributions \\
    power error (large / intermediate / small scales) & $|P_{\mathrm{pred}}/P_{\mathrm{true}}-1|$ of the transverse $\xHI$ power spectrum in three $k$ ranges; small scales mean $k>1\,\mathrm{Mpc}^{-1}$ \\
    $1-r$, high $k$ & small-scale incoherence: one minus the cross-correlation \cref{eq:cross-correlation}, averaged over $k>1\,\mathrm{Mpc}^{-1}$ \\
    edge $L^2$ / $H^1$, $\pm2$\,Mpc & error (and error in gradients) in a $\pm2$\,Mpc band around true ionization fronts \\
    front width & width of the predicted ionization fronts (Mpc), compared with the true width \\
    $T_b$ PDF $W_1$; $T_b$ power deviation & distance between $T_b$ distributions (mK); power difference in units of the single-slice sample variance ($\sigma$) \\
    \bottomrule
  \end{tabular}
\end{table*}

The cross-correlation coefficient between prediction and truth,
\begin{equation}
  \label{eq:cross-correlation}
  r(k) = \frac{P_{\mathrm{pred}\times\mathrm{true}}(k)}
              {\sqrt{P_{\mathrm{pred}}(k)\,P_{\mathrm{true}}(k)}},
\end{equation}
deserves special attention. The power ratio measures whether a model produces the right \emph{amount} of structure at each scale; $r(k)$ measures whether that structure is in the right \emph{place}. In the earlier 3-D operator study of the \texttt{fno-21cm} pipeline, small-scale coherence predicted the bias in the bubble-size distribution across architectures (Spearman $\rho=0.88$), while small-scale power amplitude did not. Getting fronts in the right place matters more for bubble statistics than getting the right amount of small-scale detail.

% ---------------------------------------------------------------------
\subsection{Baselines}
\label{sec:baselines}

All baselines come from the \texttt{fno-21cm} pipeline and were trained there on the same splits:
\begin{itemize}
  \item \textbf{\texttt{cnn\_whno\_es}} (3.5\,M parameters): the \texttt{fno} 3-D $\xHI$ model trained on native LOS windows. It has the U-FNO layout, a convolutional local path next to a learned global path, at a size comparable to the rhizome. Trained for 30 epochs (44.7\,h on an A100).
  \item \textbf{\texttt{es+histtb}} and \textbf{\texttt{tbs+histtb}}: the corresponding multi-field models, which receive density, velocity and both emulated global histories. \texttt{tbs+histtb} additionally uses the structured $T_b$ head and is the strongest \texttt{fno} model. It corresponds to rhizome variant C below and is labelled ``UFNO'' in the head-to-head figures.
  \item \textbf{2-D \texttt{fno} models} (LWF, WHNO, LocalFNO, U-FNO), which predict one slice from that slice alone.
\end{itemize}
The 203\,M-parameter implementation of U-FNO \citep{wen2022ufno} used in the earlier operator study (\cref{sec:ufno}) is not used as a baseline here, because its capacity is not comparable.

% ---------------------------------------------------------------------
\subsection{2-D slices: line-of-sight context dominates}
\label{sec:results-2d}

A sweep around a centre configuration (width 64, 24 modes, six updates) showed that the largest single factor is not a hyperparameter of the operator but the LOS context in the input (\cref{fig:los-context}). With only the target slice, the rhizome reaches a mixed-slice RMSE of 0.174. Adding the six neighbouring slices brings it to 0.138, and adding the band means out to $\pm32$ cells to 0.112, a 36\% reduction in total. By comparison, the other settings matter little: width 128 is 3\% better than width 64 at four times the parameters, 12 or 48 modes are worse than 24, and three or ten updates are slightly worse than six in this short sweep. In full-length training of the multi-field model, 3 and 10 updates match 6 within the seed-to-seed spread (\cref{sec:results-variations}). At equal compute, more training slices beat more width.

\begin{figure}[!htbp]
  \centering
  \figplaceholder[\linewidth]{los_context.pdf}{LOS context chart.}
  \caption{Effect of line-of-sight context on the 2-D rhizome (sweep, 15\,k steps, one seed). Mixed slices have mean $\xHI$ between 0.05 and 0.95.}
  \label{fig:los-context}
\end{figure}

\Cref{tab:results-2d} compares the 2-D rhizome with the 2-D \texttt{fno} models on the same 7,036 validation slices. The rhizome reduces the error of the next best model by about 39\%, and its ionization fronts are five times sharper. Most of this gap, however, is LOS information: the \texttt{fno} models see a single slice, and the single-slice rhizome (0.174) performs at their level. The 2-D comparison therefore shows what the input buys rather than what the operator buys. The fully trained 2-D model (width 128, plateau-stopped at 68\,k steps) reached a test mixed-slice RMSE of 0.099 and an ionized-region IoU of 0.914.

\begin{table}[!htbp]
  \centering
  \caption{2-D rhizome against 2-D \texttt{fno} models on 7,036 validation slices (rhizome: sweep run at 15\,k steps). Only the rhizome sees LOS context. RMSE: slice RMSE; IoU: ionized-region intersection over union; front: front width (truth 3.0\,Mpc).}
  \label{tab:results-2d}
  \small
  \begin{tabular}{@{}lrrr@{}}
    \toprule
    Model & RMSE & IoU & Front (Mpc) \\
    \midrule
    \textbf{Rhizome} (LOS bands) & \textbf{0.109} & \textbf{0.873} & \textbf{6.1} \\
    \texttt{fno} LWF & 0.178 & 0.755 & 32.6 \\
    \texttt{fno} WHNO & 0.181 & 0.755 & 31.0 \\
    \texttt{fno} LocalFNO & 0.184 & 0.742 & -- \\
    \texttt{fno} U-FNO & 0.195 & 0.722 & -- \\
    \bottomrule
  \end{tabular}
\end{table}

% ---------------------------------------------------------------------
\subsection{3-D windows: \texorpdfstring{$\xHI$}{xHI}}
\label{sec:results-3d}

The 3-D rhizome sees the same information as the 3-D \texttt{fno} model: native LOS windows of the density field. This removes the LOS-input difference, although capacity, loss and optimization still differ, so it does not isolate the interaction law alone. On the 200 test cones (\cref{tab:results-3d}), the rhizome at width 48 reaches an RMSE of 0.0594 after the warm restart, against 0.0684 for \texttt{cnn\_whno\_es}, with 5.8\,M against 3.5\,M parameters. Narrower was better (width 64 gave 0.0628), and a learning rate of $3\times10^{-4}$ was 13\% worse than $2\times10^{-3}$.

\begin{table*}[!tp]
  \centering
  \caption{3-D $\xHI$ on 200 test cones. \emph{Left:} test RMSE. \emph{Right:} physical suite (medians over cones), rhizome w48 after the warm restart.}
  \label{tab:results-3d}
  \small
  \begin{tabular}[t]{lrr}
    \toprule
    Model & Params & RMSE \\
    \midrule
    rhizome w48, 20 ep. & 5.8\,M & 0.0619 \\
    \quad + warm restart & 5.8\,M & \textbf{0.0594} \\
    rhizome w64, + restart & -- & 0.0628 \\
    \texttt{cnn\_whno\_es}, 30 ep. & 3.5\,M & 0.0684 \\
    \bottomrule
  \end{tabular}
  \hspace{1.2em}
  \begin{tabular}[t]{lrr}
    \toprule
    Metric & Rhizome & \texttt{fno} \\
    \midrule
    $|\Delta\tau_e|/\sigma_{\mathrm{Planck}}$ & \textbf{0.019} & 0.058 \\
    $|\Delta z(\bar x=0.5)|$ & \textbf{0.028} & 0.098 \\
    $b_1$ pull & \textbf{1.58} & 3.41 \\
    $P_\varepsilon$ ratio (ideal 1) & \textbf{0.89} & 0.77 \\
    bubble-size $W_1$ (Mpc) & \textbf{0.81} & 1.39 \\
    small-scale power error & \textbf{0.222} & 0.336 \\
    edge $H^1$, $\pm2$\,Mpc & \textbf{0.137} & 0.145 \\
    front width (Mpc) & \textbf{5.9} & 7.5 \\
    \bottomrule
  \end{tabular}
\end{table*}

The margin is much larger on the physical metrics than on RMSE. The rhizome is three times more accurate in $\tau_e$ and in the reionization midpoint, halves the bias error, and comes 40\% closer in bubble sizes. Its small-scale power and coherence also hold up better at every stage of reionization.

% ---------------------------------------------------------------------
\subsection{3-D multi-field: inputs, heads and losses}
\label{sec:results-3d-mf}

The multi-field models predict $\xHI$ and $T_b$ from density and LOS velocity. Four 3-D rhizome configurations were trained to separate the contributions of inputs and output head (\cref{tab:results-3d-mf}).

\begin{table}[!htbp]
  \centering
  \caption{3-D multi-field ablation, test RMSE on 199 cones ($\xHI$ / $T_b$ in mK). History channels are the band means of the emulated global $\xHI(z)$ and $T_b(z)$.}
  \label{tab:results-3d-mf}
  \footnotesize
  \setlength{\tabcolsep}{3pt}
  \begin{tabular}{@{}lp{2.6cm}rr@{}}
    \toprule
    Run & Inputs / head & 20 epochs & + restart \\
    \midrule
    A0 & density $\to\xHI$ only & 0.0602 & 0.0567 \\
    A & + velocity, + $T_b$ (plain head) & 0.0648 / 9.49 & 0.0622 / 8.71 \\
    B & A + history channels & 0.0509 / 6.26 & 0.0494 / 6.07 \\
    \textbf{C} & B + structured $T_b$ head & 0.0507 / 5.97 & \textbf{0.0495 / 5.79} \\
    \midrule
    \texttt{fno} \texttt{es+histtb} & = B & 0.0577 / 6.63 & \\
    \texttt{fno} \texttt{tbs+histtb} & = C & 0.0560 / 6.54 & \\
    \bottomrule
  \end{tabular}
\end{table}

\begin{itemize}
  \item \textbf{History channels give the largest improvement in this ablation.} Going from A to B reduces the $\xHI$ error by 21\% and the $T_b$ error by 34\%.
  \item \textbf{Velocity and a $T_b$ target alone do not help $\xHI$}; A is even slightly worse than A0 (one seed each).
  \item \textbf{The structured head mainly helps $T_b$}: without it, $T_b$ power is about 10\% low at every stage. It also improves the reionization timing by 15--30\%.
  \item \textbf{The density-only rhizome A0 already beats \texttt{fno}~\texttt{es+histtb}} (0.0567 vs 0.0577), although the latter also receives velocity and both histories.
\end{itemize}
On the physical suite, model C at 20 epochs, before the warm restart, beats \texttt{tbs+histtb} by a factor of 5 in $\tau_e$ (0.014 vs 0.076\,$\sigma$), and is closer in $b_1$ (pull 1.52 vs 2.28), $P_\varepsilon$ (0.935 vs 0.793), bubble sizes ($W_1$ 0.71 vs 1.21\,Mpc) and front width (4.53 vs 6.17\,Mpc; truth 2.77). Before reionization, \texttt{fno} adds spurious small-scale $T_b$ structure ($+12.8\sigma$ against $+3.9\sigma$). After the warm restart, model C improves further; that checkpoint is used in \cref{tab:main-results} ($\tau_e$ 0.012\,$\sigma$, $b_1$ pull 1.35, $P_\varepsilon$ 0.945, $W_1$ 0.69\,Mpc, front width 4.38\,Mpc).

\paragraph{Influence of the loss on the performance gap.} Both model families receive the same emulated global history, yet the global $\bar x_{\mathrm{HI}}(z)$ of the \texttt{fno} models sits 0.002--0.003 too low from $z\approx8$ to 25. In truly neutral cells at $z>15$ the mean error is $-0.0020$ and $-0.0030$ for the two \texttt{fno} runs against $-0.00006$ for the rhizome. A plausible contributor is the loss. For a neutral-fraction logit $\ell$, prediction $p=\mathrm{sigmoid}(\ell)$ and target $y$, $\partial_\ell\mathcal L_{\mathrm{BCE}}=p-y$, whereas a squared-error loss gives $2(p-y)p(1-p)$. Both approach zero at a correct saturated prediction, but squared error has an additional sigmoid factor that slows correction near the endpoints. This is consistent with the observed neutral-cell offset; a loss-matched retraining is needed to establish its contribution. This uniform offset integrates into a $\tau_e$ bias. A fair comparison of $\tau_e$ and the global history therefore needs the \texttt{fno} models retrained with BCE. The offset alone does not establish an explanation for the morphology differences; the effect of changing the loss on those metrics also needs to be measured.

% ---------------------------------------------------------------------
\subsection{Slice-wise reconstruction and native 3-D models}
\label{sec:results-slicewise}

The slice-wise variant (\cref{fig:slicewise}) applies the 2-D model to every slice of a cone. For a fair comparison it was retrained on a new slice cache built from the same 1,600 training cones as the 3-D models, with otherwise unchanged settings. \Cref{tab:results-slicewise} compares it with the 3-D rhizome and with \texttt{fno} on all 200 test cones.

\begin{table*}[!tp]
  \centering
  \caption{Slice-wise 2-D reconstruction against the 3-D models, $\xHI$, 200 test cones, all trained on the same 1,600 cones.}
  \label{tab:results-slicewise}
  \small
  \begin{tabular}{@{}lrrrr@{}}
    \toprule
    Metric & 2-D w128 & \textbf{2-D w64} & 3-D rhizome w48 & \texttt{cnn\_whno\_es} \\
    \midrule
    $|\Delta\tau_e|/\sigma_{\mathrm{Planck}}$ & \textbf{0.011} & \textbf{0.011} & 0.019 & 0.058 \\
    $|\Delta z(\bar x=0.5)|$ & \textbf{0.011} & 0.014 & 0.028 & 0.098 \\
    $\xHI$ PDF $W_1$, active & \textbf{0.011} & 0.013 & 0.021 & 0.030 \\
    $b_1$ pull & 1.01 & \textbf{0.85} & 1.58 & 3.41 \\
    $P_\varepsilon$ ratio (ideal 1) & \textbf{0.95} & 0.94 & 0.89 & 0.77 \\
    bubble-size $W_1$ (Mpc) & 0.59 & \textbf{0.56} & 0.81 & 1.39 \\
    power error, mid $k$ & \textbf{0.062} & 0.073 & 0.174 & 0.172 \\
    $1-r$, high $k$ & \textbf{0.119} & 0.120 & 0.199 & 0.181 \\
    edge $H^1$, $\pm2$\,Mpc & \textbf{0.109} & 0.110 & 0.138 & 0.145 \\
    front width (Mpc; truth 2.65) & \textbf{4.03} & 4.10 & 5.91 & 7.54 \\
    \bottomrule
  \end{tabular}
\end{table*}

The slice-wise model wins on nearly every physical metric, trains about seven times faster than the 3-D rhizome, and has no jitter between slices: the roughness from slice to slice matches the truth to 2--3\%, while both 3-D models are slightly over-smoothed along the LOS. Its advantage comes from front \emph{placement}: small-scale incoherence drops from 0.20 to 0.12, the mid-scale power error by 60\%, and the edge error by 21\% (significant in a paired bootstrap over cones: 2,000 resamples of the per-cone differences, 95\% percentile interval). The edge metrics use 175 of the 200 test cones: in the other 25, which reionize late, no cell falls below $\xHI=0.5$ before $z=5$, so there is no front to measure. Width barely matters, so width 64 (2.8\,h of training against 5.0\,h) is the choice. Per cone, the 2-D model wins 11 of 13 inspected cones. It loses on the earliest reionizer, cone 27 ($z_{\mathrm{mid}}\approx15.9$), where the longer context of the 3-D model helps slightly.

The error maps in \cref{fig:slicewise-errors} show the qualitative difference. The slice-wise errors are thin rims along bubble walls, whereas the 3-D models, \texttt{fno} strongly and the 3-D rhizome weakly, miss whole ionized regions, and \texttt{fno} leaves a faint ionized haze over neutral gas.

\begin{figure}[!htbp]
  \centering
  \figplaceholder[0.9\linewidth]{rhizome/slicewise_xhi_error_maps.png}{Slice-wise error maps.}
  \caption{Signed $\xHI$ error at $\bar x_{\mathrm{HI}}\approx0.4$ on eight test cones (red: predicted too neutral, blue: too ionized). Columns: truth, slice-wise 2-D rhizome (w64), 3-D rhizome (w48, warm restart), \texttt{fno} \texttt{cnn\_whno\_es}; per-panel RMSE in the corner. The bottom row is cone 27, the earliest reionizer, where all models predict ionized gas as too neutral.}
  \label{fig:slicewise-errors}
\end{figure}

For this simulation suite, a 2-D model with $\pm32$ cells ($\approx\pm46$\,Mpc) of LOS context is more accurate and cheaper than the tested 3-D windows. This does not establish a universal physical interaction length. The periodic tiling, compressed context channels, padding restrictions and optimization budget may all contribute to the result.

% ---------------------------------------------------------------------
\subsection{Multi-field slice-wise results}
\label{sec:main-results}

The final model combines the slice-wise approach with the multi-field inputs and the structured head of model C: 13 LOS bands each of density and velocity, global-history band means as scalars, and $\xHI$ and $T_b$ as outputs (\cref{sec:rhizome-variants}). It was trained on all 6,169 usable training cones (3.5\,h on an A100) and evaluated on the 199 multi-field test cones. \Cref{tab:main-results} compares it with 3-D rhizome C and with the best \texttt{fno} model.

\begin{table*}[!tp]
  \centering
  \caption{Main results, 199 multi-field test cones (medians). ``Slice-wise'' is the multi-field slice-wise 2-D rhizome trained on all cones; ``+ $T_b$ power'' adds the power term of \cref{sec:results-absorption}. 3-D rhizome C is the checkpoint after the warm restart. Best value per row in bold.}
  \label{tab:main-results}
  \small
  \setlength{\tabcolsep}{4pt}
  \begin{tabular}{@{}lrrrr@{}}
    \toprule
    Metric & Slice-wise & + $T_b$ power & 3-D rhizome C & \texttt{fno} \texttt{tbs+histtb} \\
    \midrule
    $|\Delta\tau_e|/\sigma_{\mathrm{Planck}}$ & 0.010 & \textbf{0.008} & 0.012 & 0.076 \\
    $|\Delta z(\bar x=0.5)|$ / duration & \textbf{0.010 / 0.018} & \textbf{0.010 / 0.018} & 0.011 / \textbf{0.018} & 0.017 / 0.044 \\
    $\xHI$ PDF $W_1$, active & \textbf{0.0095} & 0.0098 & 0.0156 & 0.0271 \\
    $b_1$ pull & 0.99 & \textbf{0.93} & 1.35 & 2.28 \\
    $P_\varepsilon$ ratio (ideal 1) & \textbf{0.954} & -- & 0.945 & 0.793 \\
    bubble-size $W_1$, active (Mpc) & 0.56 & \textbf{0.48} & 0.69 & 1.21 \\
    $\xHI$ power error, high $k$ & 0.168 & \textbf{0.163} & 0.223 & 0.311 \\
    $1-r$, high $k$ & \textbf{0.096} & 0.102 & 0.151 & 0.174 \\
    edge $H^1$, $\pm2$\,Mpc & \textbf{0.106} & 0.112 & 0.131 & 0.147 \\
    front width (Mpc; truth 2.77) & \textbf{3.89} & 3.94 & 4.38 & 6.17 \\
    $T_b$ PDF $W_1$, active (mK) & 0.79 & \textbf{0.77} & 1.06 & 1.19 \\
    pre-reion.\ $T_b$ power dev., $k=0.46$ & $+24.8\sigma$ & $\bm{+1.5\sigma}$ & $+4.1\sigma$ & $+12.8\sigma$ \\
    \bottomrule
  \end{tabular}
\end{table*}

Against \texttt{fno}, the slice-wise model reduces the $\xHI$ value-distribution error by 65\%, the $b_1$ pull by 57\%, the bubble-size error by 54\% and small-scale incoherence by 45\%, and its fronts are 3.9\,Mpc wide against 6.2\,Mpc. During reionization, its $T_b$ distribution is 34\% closer to the truth. On the eight test cones inspected individually, it beats \texttt{fno} on every cone, by 12--28\% in full-cone RMSE. It also beats the 3-D rhizome C on every $\xHI$ morphology metric (value distribution $-39\%$, $b_1$ pull $-27\%$, bubble sizes $-19\%$, incoherence $-36\%$, edges $-19\%$), while $\tau_e$ and timing tie: with history channels, both rhizome models are already five to seven times better than \texttt{fno} there. Training on all cones instead of the 1,594-cone split improved per-cone RMSE by 3.6\% ($\xHI$) and 5.5\% ($T_b$) in the median, on 86\% and 75\% of cones respectively.

\begin{figure*}[!tp]
  \centering
  \figplaceholder[0.56\linewidth]{rhizome/h2h_cone24.png}{Cone 24 head-to-head.}
  \caption{Test cone 24. \emph{Top left:} $\xHI$ lightcone strips for truth, rhizome (slice-wise multi-field) and the \texttt{fno} baseline (labelled UFNO). \emph{Top right:} global history and signed error strips. \emph{Bottom:} transverse slices at mean neutral fraction 0.17, 0.40, 0.61 and 0.80, with per-slice RMSE. Full-cone RMSE 0.045 (rhizome) against 0.057 (\texttt{fno}). The \texttt{fno} fronts are blurred, and it leaves a faint haze inside ionized regions.}
  \label{fig:cone24}
\end{figure*}

\Cref{fig:h2h-ps} shows the scale dependence. Both models lose small-scale power late in reionization, \texttt{fno} more so. The rhizome keeps $r>0.9$ to higher $k$ at every stage, so the structure it does produce sits in the right place. \Cref{fig:ps-plane} puts this in the context of all models from the earlier 3-D operator study. In the plane of small-scale coherence against small-scale power error, the multi-field slice-wise rhizome sits apart from every \texttt{fno} model: $r(k>1)=0.90$ against 0.82--0.83, power error 0.17 against 0.31--0.34.

\begin{figure*}[!tp]
  \centering
  \figplaceholder[\linewidth]{rhizome/h2h_ps_stage_curves.png}{Power and coherence by stage.}
  \caption{$\xHI$ power ratio (left) and cross-correlation $r(k_\perp)$ (right) by reionization stage, multi-field slice-wise rhizome (top) and \texttt{fno} \texttt{tbs+histtb} (bottom), 199 test cones.}
  \label{fig:h2h-ps}
\end{figure*}

\begin{figure*}[!tp]
  \centering
  \figplaceholder[\linewidth]{rhizome/ps_plane_coherence_vs_power_best.png}{Coherence vs power plane.}
  \caption{Small-scale coherence $r(k>1\,\mathrm{Mpc}^{-1})$ against small-scale power error for the 3-D $\xHI$ models of the earlier operator study (open circles), the \texttt{fno} LOS-window models (filled) and the multi-field slice-wise rhizome (green), active stage. Better models lie to the lower right. Models from different pipelines and splits are compared with care.}
  \label{fig:ps-plane}
\end{figure*}

\begin{figure*}[!tp]
  \centering
  \figplaceholder[\linewidth]{rhizome/h2h_physical_history.png}{History and timing.}\\[4pt]
  \figplaceholder[\linewidth]{rhizome/h2h_bubble_size_summary.png}{Bubble size summary.}\\[4pt]
  \figplaceholder[\linewidth]{rhizome/h2h_edge_3d.png}{Edge profiles.}
  \caption{Physical metrics, rhizome (slice-wise multi-field) against \texttt{fno} \texttt{tbs+histtb}, 199 test cones. \emph{Top:} global history error, $\tau_e$ error in Planck $\sigma$ (median 0.01 against 0.08), and predicted against true midpoint and duration. \emph{Middle:} bubble-size distributions and their $W_1$ distance to the truth by stage (active 0.56 against 1.21\,Mpc). \emph{Bottom:} error profiles across ionization fronts; front width 3.9 against 6.2\,Mpc (truth 2.77).}
  \label{fig:h2h-physical}
\end{figure*}

\Cref{fig:h2h-errors} shows the signed error at $\bar x_{\mathrm{HI}}\approx0.4$ on eight test cones. As for the $\xHI$-only model, the rhizome's errors are thin rims on bubble walls, while \texttt{fno} misses whole regions and its rims are wider.

\begin{figure}[!htbp]
  \centering
  \figplaceholder[0.78\linewidth]{rhizome/h2h_error_maps_xhi04.png}{Error maps.}
  \caption{Signed $\xHI$ error maps at $\bar x_{\mathrm{HI}}\approx0.4$ on eight test cones (red: predicted too neutral), rhizome against \texttt{fno} \texttt{tbs+histtb}.}
  \label{fig:h2h-errors}
\end{figure}

\paragraph{Shared failure.} Some test cones contain large, flat plateaus of intermediate neutral fraction ($\xHI\approx0.3$--0.6) in the truth. No model reproduces them, and all overpredict the partially ionized fraction (0.13--0.18 against 0.09).

% ---------------------------------------------------------------------
\subsection{Brightness temperature before reionization}
\label{sec:results-absorption}

The slice-wise model has one clear weakness. Before reionization, when $T_b$ is in absorption, it produces too much small-scale $T_b$ power: $+24.8\sigma$ at $k=0.46\,\mathrm{Mpc}^{-1}$, against $+4.1\sigma$ for the 3-D rhizome and $+12.8\sigma$ for \texttt{fno} (\cref{tab:main-results}). The excess is concentrated in quiet slices whose true small-scale power is tiny; the power summed over the cone is in fact slightly low (ratio 0.97). The working explanation is that the spin variable $u$ is predicted slice by slice, and the per-pixel loss has no explicit term matching its fluctuation spectrum within a slice, whereas the 3-D model predicts it jointly along the LOS.

Three fixes that address data, post-processing and context did not work; a fourth, in the loss, did (\cref{tab:absorption-fixes}).

\begin{table}[!htbp]
  \centering
  \caption{Attempts to remove the absorption-era $T_b$ power excess of the slice-wise model (power deviation at $k=0.46\,\mathrm{Mpc}^{-1}$, plain model: $+24.8\sigma$).}
  \label{tab:absorption-fixes}
  \footnotesize
  \begin{tabular}{@{}p{2.2cm}p{1.05cm}p{4.3cm}@{}}
    \toprule
    Attempt & Excess & Outcome \\
    \midrule
    3.9$\times$ more training cones & $+24.8\sigma$ & $T_b$ values improve ($W_1$ $-31\%$), excess unchanged: no improvement from this data increase \\
    smooth $u$ along the LOS ($\sigma=4$ cells) & $+1.1\sigma$ & removes the excess only by destroying real LOS structure: $T_b$ RMSE $+40$--$110\%$, summed power $\times1.3$; rejected \\
    LOS context $\pm128$ cells & $+36.0\sigma$ & worse; per-voxel accuracy unchanged; not adopted \\
    \textbf{small-scale $T_b$ power loss} & $\bm{+1.5\sigma}$ & reduces the reported excess across scales; costs 4--8\% per-voxel RMSE \\
    \bottomrule
  \end{tabular}
\end{table}

The power loss compares the transverse power spectrum of each predicted slice with the truth in eight logarithmic bins from $k=0.3\,\mathrm{Mpc}^{-1}$ to the grid Nyquist frequency,
\begin{equation}
  \label{eq:power-loss}
  \mathcal{L}_{P} = \lambda\,\Big\langle \big(\log(P_{\mathrm{pred}}+\epsilon)-\log(P_{\mathrm{true}}+\epsilon)\big)^2 \Big\rangle_{\text{slices, bins}},
\end{equation}
with $\lambda=0.03$, applied to the mean-subtracted normalized $T_b$. It brings the absorption-era excess from $+24.8\sigma$ to $+1.5\sigma$, better than the 3-D rhizome, without overshooting the amplitude. It also improves the reionization-era statistics: $T_b$ power within 1--2\% at all scales, bubble-size $W_1$ $-14\%$ and $\tau_e$ $-20\%$ (\cref{fig:power-loss}). The price is a 4--8\% higher per-voxel RMSE, consistent over 80--90\% of the cones, and slightly worse edges. This reflects the different quantities optimized by the two losses: squared error rewards agreement at each cell, while the power term rewards agreement in fluctuation amplitude. Matching a power spectrum does not guarantee correct phases or morphology, so cross-correlation and front metrics remain necessary.

Applying the same term to $\xHI$ as well roughly halves the high-$k$ $\xHI$ power deficit (from 0.163 to 0.089 over active slices; in the latest stage from 0.59 to 0.15), gives the closest value distribution of all models ($W_1$ 0.0082), and keeps the absorption-era fix ($+1.3\sigma$). It costs another 2--4\% in per-voxel RMSE, and the extra small-scale power is only partly in phase with the truth (small-scale incoherence $+7\%$). The three variants are useful for different purposes: the plain model for per-voxel accuracy, the $T_b$ power loss for 21\,cm statistics, and both power terms for closer small-scale power amplitudes in both fields, with the stated phase-error trade-off.

\begin{figure*}[!tp]
  \centering
  \figplaceholder[\linewidth]{rhizome/power_loss_physical_21cm.png}{Power-loss 21cm statistics.}
  \caption{21\,cm statistics of the power-loss variants during reionization, 199 test cones: plain slice-wise model, $+\,T_b$ power loss (\texttt{pw0.03}), $+\,T_b$ and $\xHI$ power loss (\texttt{pw0.03\_px0.03}), and 3-D rhizome C. \emph{Right column:} power bias in units of single-slice sample variance for $T_b$ (top) and $\xHI$ (bottom). Both power-loss models stay near zero to $k\approx2\,\mathrm{Mpc}^{-1}$ in $T_b$; the plain model and 3-D C fall to $-10$ and $-12\sigma$.}
  \label{fig:power-loss}
\end{figure*}

A further third of the absorption-era $T_b$ error is a constant offset per slice, of random sign from cone to cone. Its size (1.6--2.0\,mK RMS) is nearly the same for all three model families and comparable to the error of the global-history emulator that provides the $\bar T_b(z)$ input (about 2.1\,mK). The similar error scales suggest that the history input may limit this component. Testing with true rather than emulated histories would separate that effect from model error. The power loss acts only on fluctuations and leaves the offset unchanged.

% ---------------------------------------------------------------------
\subsection{Controlled variations of the slice-wise model}
\label{sec:results-variations}

To separate components of the final model, the plain multi-field slice-wise configuration was retrained from scratch with one change at a time: a second seed, the number of updates, the interaction kernel, and the history inputs. Data, splits, losses, stopping rule and evaluation are the same as for the reference run (\cref{tab:variations}). The variants were trained with a later code version whose additions are inactive at these settings; the seed-1 run uses that same version. The kernel variants replace the global spectral kernel by a \emph{compact} real-space kernel supported on a disc of radius $R$ cells. All gates, the recurrence and the pointwise maps are kept. After $\Niter=6$ updates, information can then travel at most $6R$ cells. The rank-16 variant factorizes the channel mixing of the compact kernel.

\begin{table*}[!tp]
  \centering
  \caption{Controlled variations of the multi-field slice-wise model (width 64, all training cones, plain loss). The reference uses the spectral kernel with 24 modes, 6 updates and seed 0; each other row changes one setting and is retrained from scratch. Test RMSE on 3,129 held-out slices from 199 test cones; $\Delta$ relative to the seed-0 reference. Reach: largest distance information can travel in six updates. $P$ error: mean squared log ratio of the transverse power spectrum, predicted to true, in eight bins above $k=0.3\,\mathrm{Mpc}^{-1}$ (not recorded for the reference run). Training time on one A100, including validation.}
  \label{tab:variations}
  \footnotesize
  \setlength{\tabcolsep}{4pt}
  \begin{tabular}{@{}lrrrrrr@{}}
    \toprule
    Variant & Params & Reach & Time (h) & $\xHI$ RMSE ($\Delta$) & $T_b$ RMSE, mK ($\Delta$) & $P$ error $T_b$ / $\xHI$ \\
    \midrule
    reference (seed 0) & 4.57\,M & global & 3.6 & 0.0872 & 7.71 & -- \\
    second seed (seed 1) & 4.57\,M & global & 3.1 & 0.0873 ($+0.2\%$) & 7.82 ($+1.3\%$) & 0.385 / 0.392 \\
    3 updates & 4.57\,M & global & 2.2 & 0.0866 ($-0.6\%$) & 7.68 ($-0.5\%$) & 0.395 / 0.385 \\
    10 updates & 4.57\,M & global & 4.9 & 0.0869 ($-0.3\%$) & 7.66 ($-0.7\%$) & 0.378 / 0.355 \\
    \midrule
    compact kernel, $R=1$ & 66\,k & 8.6\,Mpc & 2.8 & 0.1209 ($+38.7\%$) & 9.76 ($+26.5\%$) & 0.511 / 0.449 \\
    compact kernel, $R=2$ & 99\,k & 17\,Mpc & 3.8 & 0.1089 ($+24.9\%$) & 8.85 ($+14.7\%$) & 0.424 / 0.398 \\
    compact kernel, $R=3$ & 164\,k & 26\,Mpc & 4.4 & 0.1010 ($+15.9\%$) & 8.27 ($+7.3\%$) & 0.424 / 0.402 \\
    compact kernel, $R=4$ & 246\,k & 34\,Mpc & 5.2 & 0.0972 ($+11.5\%$) & 8.10 ($+5.1\%$) & 0.417 / 0.385 \\
    compact kernel, $R=6$ & 508\,k & 51\,Mpc & 7.7 & 0.0951 ($+9.1\%$) & 8.13 ($+5.4\%$) & 0.444 / 0.437 \\
    compact kernel, $R=8$, rank 16 & 51\,k & 69\,Mpc & 4.6 & 0.0926 ($+6.3\%$) & 8.04 ($+4.2\%$) & 0.444 / 0.434 \\
    \midrule
    no $\bar x_{\mathrm{HI}}(z)$ history input & 4.57\,M & global & 3.7 & 0.0973 ($+11.7\%$) & 8.26 ($+7.1\%$) & 0.399 / 0.431 \\
    no $\bar T_b(z)$ history input & 4.57\,M & global & 3.9 & 0.0900 ($+3.3\%$) & 8.04 ($+4.3\%$) & 0.341 / 0.400 \\
    \bottomrule
  \end{tabular}
\end{table*}

\paragraph{Seed and rollout length.} The second seed differs from the reference by 0.2\% in $\xHI$ and 1.3\% in $T_b$ RMSE. This is the only direct estimate of training noise in this report, and it suggests that differences of one or two percent between single runs should not be interpreted. Three and ten updates match six within that spread. Three updates train in 2.2\,h instead of 3.6\,h. Between 3 and 10 updates, the rollout length is therefore not a sensitive setting on this task. This differs from the toy pilot of \cref{tab:rhizome-pilot}, where 1 or 2 updates were clearly worse, but there each update count was evaluated without retraining.

\paragraph{Global reach of the kernel.} Every compact kernel is worse than the spectral kernel, and the error falls steadily as the reach grows: from $+39\%$ in $\xHI$ at 8.6\,Mpc to $+6\%$ at 69\,Mpc. The compact kernels also have 9--90 times fewer parameters, so reach and capacity change together. Within the compact family, however, the error follows reach rather than parameter count: the rank-16 kernel with $R=8$ has the fewest parameters (51\,k) and the smallest error, while $R=6$ with ten times more parameters is worse. The compact kernels are not cheaper to train on an A100 either, because wide direct convolutions are slower than the FFT (7.7\,h at $R=6$). These runs support the global-reach design principle of \cref{sec:rhizome-principles}. Since every compact variant keeps the gates, they do not measure the value of the gates.

\paragraph{Global-history inputs.} Both histories were removed in turn from the scalar inputs. Paired over the 199 full test cones, removing the $\xHI$ history raises the median per-cone RMSE by 12.6\% in $\xHI$ (95\% bootstrap interval 11.5--14.5\%) and by 11.3\% in $T_b$; it is worse on 91\% and 93\% of the cones. The error of the global history $\bar x_{\mathrm{HI}}(z)$ grows by 38\%. Removing the $T_b$ history raises the median per-cone $T_b$ error by 26.6\% (interval 19.7--31.2\%), and by 46\% before reionization. Its $\xHI$ error rises by 6.1\%. Both changes are far larger than the seed spread. They agree with the 3-D ablation of \cref{tab:results-3d-mf}: the emulated histories carry information that the local density bands do not. One effect goes the other way: without the $T_b$ history the slice-level $T_b$ power error is lower (0.341 against 0.378--0.395), so the history input improves per-voxel accuracy but not the small-scale fluctuation amplitude.

% ---------------------------------------------------------------------
\subsection{Learned spectral weights}
\label{sec:results-modes}

Neural networks tend to learn low frequencies first, and a Fourier layer can end up using only its lowest modes. \Cref{fig:mode-weights} tests this for the rhizome by comparing the RMS weight of each Fourier mode of the trained kernel with its value at initialization. The weights grew 13--45-fold almost evenly up to the band edge: the ratio of high-quarter to low-quarter weight is 0.78--0.94 for all rhizome models, against 0.10--0.45 for a U-FNO in an earlier study. Models with 36 or 48 modes also develop non-negligible weights in their extra modes. Weight amplitude alone, however, does not measure mode usage: the transmitted input spectrum and output sensitivity also matter. These profiles show no collapse of the multiplier weights, but do not establish why the 48-mode sweep performed worse or rule out an alternative kernel parameterization. Longer matched-budget runs and per-mode message or sensitivity measurements would test those explanations.

\begin{figure*}[!tp]
  \centering
  \figplaceholder[0.82\linewidth]{rhizome/mode_weight_profiles.png}{Mode weight profiles.}
  \caption{Per-mode RMS of the learned Fourier kernel relative to initialization. Radial profiles over wavenumber; they are flat to the band edge: there is no collapse onto low frequencies.}
  \label{fig:mode-weights}
\end{figure*}

% ---------------------------------------------------------------------
\subsection{Learned gates on 21\,cm slices}
\label{sec:results-gates}

\Cref{fig:mechanism} inspects the gates of a toy model. The same analysis was repeated on the trained multi-field slice-wise model, averaged over 120 test slices with $0.1\le\bar x_{\mathrm{HI}}\le0.9$, one per cone. Receiver cells are classified with the true field: 3,808 cells on ionization fronts, and 1,920 and 1,489 cells deep inside neutral and ionized regions. For each receiver $i$ and update $t$, the re-weighting $M_t(j)$ of \cref{fig:mechanism} is averaged radially in $|r_i-r_j|$ (\cref{fig:mf-gates}).

\begin{figure*}[!tp]
  \centering
  \figplaceholder[0.82\linewidth]{rhizome/mf_gate_gain_receiver.png}{Gate gain and receiver gate per update.}
  \caption{Gates of the trained multi-field slice-wise model on 120 test slices, by receiver class and update. \emph{Left:} mean gain $\langle M_t\rangle_j$ of the incoming connections. \emph{Right:} channel-RMS receiver gate $\norm{A_i}_{\mathrm{rms}}$. Solid: plain model; dashed: model trained with both power terms (\cref{sec:results-absorption}). Gains on the transmitted value, not output sensitivities.}
  \label{fig:mf-gates}
\end{figure*}

\paragraph{The gates depend on the state.} The receiver gate separates the phases, and the separation grows during the rollout. In the plain model, the channel-RMS receiver gate of ionized-interior cells rises from 0.34 to 0.46 over the six updates, while that of neutral-interior cells falls from 0.26 to 0.15 and that of front cells from 0.25 to 0.17. At the last update ionized cells accept about three times more of the incoming message than neutral ones. The source gates stay close to one (channel RMS 0.81--0.97 on four inspected slices), so, as in the toy model, the receiver dominates. The gates do not merely rescale: the radial profile of the re-weighting also changes. For front receivers, the boosted region contracts from a half-width of 26\,Mpc at the first update to 14\,Mpc at the last, with a peak boost of $2.0\times$ at 3.6\,Mpc; for neutral-interior receivers it contracts from 28 to 5\,Mpc. For ionized-interior receivers it stays broad (25 to 24\,Mpc). Beyond 40\,Mpc the re-weighting is flat for all classes. On four single slices, the normalized shape of the effective connection map into a front cell changes by 25--30\% between the first and last update, against about 4\% in the toy model. The model trained with both power terms has larger receiver gates and concentrates its re-weighting more strongly, within about 1\,Mpc of the receiver ($3.5\times$ at fronts).

\paragraph{What this does and does not show.} The trained gates are clearly state dependent: they distinguish ionized from neutral cells, and the connections they produce change over the rollout, mostly by narrowing towards nearby sources in neutral and front cells. This is a description of the learned mechanism, not a measurement of its value. A model with gates computed once from the initial state, or without gates, could reach the same accuracy by other means, and $M_t$ measures gains on the transmitted value rather than influence on the output.

% ---------------------------------------------------------------------
\subsection{Cost}
\label{sec:results-cost}

Per training slice, the 2-D model is about four times slower than the 3-D models, since it runs full-resolution recurrent updates in float32. It needs 15--20 times fewer training slices, however (2.2\,M against 33--49\,M), so a complete training run is four to seven times cheaper: 3.5\,h on one A100 for the best model, against about 18\,h on an H200 for the 3-D rhizome and 44.7\,h on an A100 for \texttt{cnn\_whno\_es}.

At inference, a full multi-field cone takes 5.5\,s end to end on an A100 and 3.7\,s on an H200 once the LOS bands are built on the GPU and the next data window is prefetched (\cref{fig:inference-cost}). This makes the slice-wise model the fastest model end to end, 1.8--2.3 times faster than both 3-D models; at that point inference is limited by reading the data. The GPU forward pass alone (3.8\,s / 1.7\,s) matches the 3-D rhizome and is about twice as slow as \texttt{fno}. Mixed precision is not worthwhile: bfloat16 gains only 5--14\%, because the FFTs stay in float32, and it shifts $\xHI$ by up to 0.08.

\begin{figure}[!htbp]
  \centering
  \figplaceholder[\linewidth]{inference_cost.pdf}{Inference cost chart.}
  \caption{End-to-end inference time for one full multi-field test cone ($\xHI$ and $T_b$), including data reading.}
  \label{fig:inference-cost}
\end{figure}

% ---------------------------------------------------------------------
\subsection{Ablations}
\label{sec:ablations}

\Cref{tab:ablations} collects what each design and training choice was worth across the experiments above. The table includes controlled changes as well as comparisons of model configurations. All training runs use a single seed, except one seed repeat of the main configuration, which differs by 0.2\% ($\xHI$) and 1.3\% ($T_b$) in test RMSE. The 2-D/3-D comparison changes several architectural and numerical choices together.

\begin{table*}[!tp]
  \centering
  \caption{Summary of ablations. Effects are relative changes in the stated metric; negative is better.}
  \label{tab:ablations}
  \small
  \begin{tabular}{@{}p{4.6cm}p{3.6cm}p{5.6cm}@{}}
    \toprule
    Change & Metric & Effect \\
    \midrule
    LOS bands vs single slice (2-D) & mixed-slice RMSE & $-36\%$ (0.174 $\to$ 0.112) \\
    width 64 $\to$ 128 (2-D) & mixed-slice RMSE & $-3\%$ at $4\times$ parameters \\
    width 48 $\to$ 64 (3-D) & test RMSE & $+6\%$ (worse) \\
    lr $2\times10^{-3}\to3\times10^{-4}$ (3-D) & test RMSE & $+13\%$ (worse) \\
    warm restart (3-D) & test RMSE & $-3$ to $-6\%$ \\
    history channels (3-D, A $\to$ B) & RMSE $\xHI$ / $T_b$ & $-21\%$ / $-34\%$ \\
    structured $T_b$ head (B $\to$ C) & $T_b$ RMSE; timing & $-5\%$; $-15$ to $-30\%$ \\
    slice-wise 2-D vs 3-D (same data) & $1-r$; edge $H^1$ & $-40\%$; $-21\%$ \\
    all 6,169 vs 1,594 training cones & median per-cone RMSE & $-3.6\%$ ($\xHI$), $-5.5\%$ ($T_b$) \\
    LOS context $\pm128$ vs $\pm32$ & per-voxel RMSE & little reported change; duration $-37\%$; absorption $T_b$ worse \\
    $T_b$ power loss & pre-reion.\ $T_b$ power dev. & $+24.8\sigma\to+1.5\sigma$; RMSE $+4$ to $+8\%$ \\
    $+\,\xHI$ power loss & high-$k$ $\xHI$ power error & $-45\%$; RMSE further $+2$ to $+4\%$ \\
    \midrule
    \multicolumn{3}{@{}l}{\emph{Slice-wise multi-field model, retrained (\cref{sec:results-variations})}} \\
    second seed & test RMSE $\xHI$ / $T_b$ & $+0.2\%$ / $+1.3\%$ (seed spread) \\
    3 or 10 updates instead of 6 & test RMSE $\xHI$ / $T_b$ & $-0.3$ to $-0.6\%$ / $-0.5$ to $-0.7\%$ (within seed spread) \\
    compact kernel, reach 8.6--69\,Mpc, vs spectral & test RMSE $\xHI$ / $T_b$ & $+39\%$ to $+6\%$ / $+27\%$ to $+4\%$ (worse; shrinks with reach) \\
    remove $\bar x_{\mathrm{HI}}(z)$ history & median per-cone RMSE $\xHI$ / $T_b$ & $+12.6\%$ / $+11.3\%$ (worse) \\
    remove $\bar T_b(z)$ history & median per-cone RMSE $\xHI$ / $T_b$ & $+6.1\%$ / $+26.6\%$ (worse) \\
    \bottomrule
  \end{tabular}
\end{table*}

\paragraph{What the variations establish about the operator.} The retrained variations of \cref{sec:results-variations} isolate three components of the operator. The global reach of the kernel matters: no compact kernel up to 69\,Mpc reach matches the spectral kernel. The rollout length matters little between 3 and 10 updates. The history inputs carry information the operator cannot recover from the density bands. \Cref{sec:results-gates} shows that the trained gates are state dependent and change over the rollout.

\paragraph{What remains open.} None of these runs removes or freezes the gates, so the accuracy contribution of the state-dependent gating is still not measured. The planned study compares the full rhizome with six otherwise identical arms (no gates, source-only, receiver-only, gates computed from the conditioning only, gates computed once from the initial state, and static per-channel gates), five seeds each, all retrained from scratch. With the seed spread now measured at 1--2\%, the primary contrast, full against initial-state gates, needs differences well above that level to be interpretable. Until then, the comparison with \texttt{fno} establishes that the rhizome \emph{as a whole} works better on this task, and the variations attribute part of that to the global kernel and the inputs, but not to the state-dependent gating specifically.

% ---------------------------------------------------------------------
\subsection{Discussion}
\label{sec:discussion}

\paragraph{What the rhizome does well.} Across every configuration, the rhizome's advantage over the \texttt{fno} models is larger on physical observables than on per-voxel RMSE. It places ionization fronts more accurately (higher small-scale coherence), keeps them sharper, and reproduces bubble-size distributions, bias and stochasticity more faithfully. These are among the quantities that \cref{sec:fno-limitations} identified as useful diagnostics of a restricted spectral interaction. They are also the ones that matter for 21\,cm inference, where bubble morphology and the field-level bias enter the observables.

\paragraph{What the line of sight needs.} The most consequential finding is not about the operator itself. A 2-D model with $\pm32$ cells of LOS context, applied slice by slice, outperforms native 3-D models trained on the same data, and does so at a fraction of the cost. Together with the periodic tiling of the lightcones, this suggests that the bottleneck of the 3-D models is how they handle the LOS, not their transverse expressivity.

\paragraph{Scope of the evidence.} The targets are semi-numerical lightcones from one simulation suite with repeated periodic boxes. Accuracy on those targets does not by itself establish agreement with resolved radiative transfer, unseen cosmologies, another box size or observational inference. The emulator is deterministic and has no calibrated field uncertainty. Power spectra, morphology diagnostics and optical-depth errors assess different aspects of agreement, but do not replace a downstream inference test.

\paragraph{Reproducibility details still needed.} The optimizer, stopping rule, normalization, loss weights and slice sampling of the slice-wise models are now recorded in \cref{sec:training}. Before treating this report as a complete methods paper, the simulation-parameter ranges and sampling law, split identifiers, the corresponding settings of the 3-D models, history-emulator training split, and precise metric definitions should also be recorded. The history inputs need an explicit provenance check showing that test-cone targets were not used to fit or select the emulator. LOS padding length must be distinguished from core and halo lengths. The meaning of the quoted $\sigma$ and $b_1$ pulls, the partially ionized thresholds, and power-bin boundaries also need to be specified. These are missing documentation or checks, not assumptions supplied by this revision.

\TODO{Record whether the test comparisons informed the choice of context, loss weights or checkpoints.}

\paragraph{Caveats.} Several confounds remain. Nearly all results are single seeds. A repeat of the main slice-wise configuration differs by 0.2--1.3\% in test RMSE, so differences of a few percent (for example A against A0, which is a 3-D run with its own noise level) may be partly noise. Part of the advantage in $\tau_e$ and the global history comes from training with BCE rather than from the operator, and the \texttt{fno} models should be retrained with BCE to separate the two. The contribution of the gating itself has not been isolated (\cref{sec:ablations}), although the trained gates are demonstrably state dependent (\cref{sec:results-gates}). Finally, all results are at one resolution; the grid-independence of the parameters has not been put to the test.




\FloatBarrier
\appendix
% =====================================================================
\section{Notation}
\label{app:notation}
% =====================================================================

\begin{table}[H]
  \centering
  \footnotesize
  \begin{tabular}{@{}ll@{}}
    \toprule
    Symbol & Meaning \\
    \midrule
    \multicolumn{2}{@{}l}{\emph{Operator learning}} \\
    $f,\ g$                   & input and output functions \\
    $\A,\ \U$                 & input and output function spaces \\
    $D,\ N$                   & spatial domain and total grid-cell count \\
    $\kappa(r,r',\theta)$     & kernel of an integral operator \\
    $\F,\ \Finv$              & Fourier transform and inverse \\
    $R_\phi(k)$               & learned Fourier coefficients of the kernel (FNO) \\
    $\kmax,\ \Lambda$           & mode-count setting and retained index set \\
    $\dv$                     & channel width \\
    $P,\ Q,\ W,\ \sigma$      & lifting, projection, pointwise linear path, nonlinearity \\
    \midrule
    \multicolumn{2}{@{}l}{\emph{Rhizome operator}} \\
    $\h^t,\ \cond,\ \msg^t$   & hidden state at update $t$, conditioning, message \\
    $\Agate,\ \Bgate$         & receiver and source gates, values in $(0,2)$ \\
    $\Vval$                   & value map \\
    $\kernel,\ R(k)$          & base spatial kernel and its Fourier coefficients \\
    $\kappa^{\mathrm{eff}}_t$ & effective, state-dependent kernel at update $t$ \\
    $\Niter,\ \Delta t$       & computational updates (6), step size (0.5) \\
    \midrule
    \multicolumn{2}{@{}l}{\emph{21\,cm physics and evaluation}} \\
    $\xHI$                    & neutral hydrogen fraction; $\bar x_{\mathrm{HI}}$ its global mean \\
    $T_b$                     & 21\,cm brightness temperature (mK) \\
    $\delta,\ z,\ \vartheta$  & density contrast, redshift, simulation parameters \\
    $u$                       & spin variable, $1-e^{u}$ the spin-temperature factor \\
    $T_0(z),\ F_v$            & $T_b$ amplitude and velocity-gradient factor \\
    $\tau_e$                  & CMB optical depth to reionization \\
    $r(k)$                    & cross-correlation coefficient of prediction and truth \\
    $b_1,\ P_\varepsilon$     & linear bias and stochastic power of the EFT expansion \\
    $W_1$                     & Wasserstein-1 distance \\
    LOS                       & line of sight \\
    \bottomrule
  \end{tabular}
\end{table}

\FloatBarrier

{\small
\bibliographystyle{plainnat}
\bibliography{references}}

\end{document}
